% Kombinierbarer Einschub fuer ein groesseres Hauptdokument.
% Vorgesehen fuer % Kombinierbarer Einschub fuer ein groesseres Hauptdokument.
% Vorgesehen fuer % Kombinierbarer Einschub fuer ein groesseres Hauptdokument.
% Vorgesehen fuer % Kombinierbarer Einschub fuer ein groesseres Hauptdokument.
% Vorgesehen fuer \input{Abschluss.tex} oder \include{Abschluss}.

\section{Numerische Unterlegung von Section 7}
\label{sec:section7_numerische_unterlegung}

Die bisherige inhaltliche Struktur von Section~7 laesst sich nun auch
numerisch unterlegen. Ziel ist es, die wesentlichen Marker nicht nur
qualitativ zu benennen, sondern ihre Stabilitaet in zwei verschiedenen
Regimen sowie ihre Kopplung an die Nullstellenseite sichtbar zu machen.

\subsection{Erzeugte Dateien}

Die numerische Auswertung fuehrt auf vier zentrale Ausgabedateien:
\begin{itemize}
\item \texttt{Section 7 -- Bericht},
\item \texttt{Alle Extrapolationswerte},
\item \texttt{Markeruebersicht fuer $s=1$ und $s=-1$},
\item \texttt{Fenstervergleich zur Nullstellenseite}.
\end{itemize}

Diese Dateien bilden zusammen die numerische Basis fuer die folgende
Interpretation.

\subsection{Untersuchte Observablen}

Betrachtet wurden fuenf zentrale Observablen:
\begin{itemize}
\item $K_{\mathrm{add}}$,
\item $G_{P,235}$,
\item $\mathrm{Bridge}_{P,M}$,
\item $\mathrm{Bridge}_{P,W}$,
\item $d_e$ als Abstand zur naechsten E-Primstelle.
\end{itemize}

Fuer jede dieser Groessen wurden zwei komplementaere Regime parallel
untersucht.

\subsection{Das $s=1$-Regime}

Im ersten Regime wurden die Partialsummen
\[
\sum_{p\le X}\frac{F(p)}{p}
\]
gegen
\[
\log\log X+\gamma
\]
gefittet. Die Euler-Mascheroni-Konstante
\[
\gamma \approx 0.5772156649
\]
fungiert dabei nicht nur als kosmetische Verschiebung, sondern als
explizites Zentrum des harmonischen beziehungsweise polnahen
$s=1$-Verhaltens.

\subsection{Das $s=-1$-Regime}

Im zweiten Regime wurden die rohen Summen
\[
\sum_{p\le X} F(p)\,p
\]
mit einem Trendmodell der Form
\[
a\,X\log X+b\,X+c\,\log X+d
\]
approximiert. Anschliessend wurden die Residuen untersucht. Diese
Regularisierung ist zwar noch keine analytische Fortsetzung im strengen
Sinn, sie ist numerisch aber ausreichend stabil, um zwischen ruhigen und
unruhigen Kanaelen zu unterscheiden.

\subsection{Lesart der Markeruebersicht}

Die zentrale Tabelle ist die Markeruebersicht fuer $s=1$ und $s=-1$.
Sie enthaelt fuer jede Observable insbesondere:
\begin{itemize}
\item die Steigung im $\gamma$-zentrierten $s=1$-Fit,
\item den letzten Residualwert,
\item die Gesamtbreite der Residuen,
\item die Koeffizienten des Trendmodells im $s=-1$-Regime.
\end{itemize}

Praktisch ist die Interpretation einfach:
\begin{itemize}
\item Eine kleine \texttt{s1\_residual\_range} bedeutet, dass sich der
Kanal gut an das Bild $\log\log X+\gamma$ anschliesst.
\item Eine kleine \texttt{sminus1\_residual\_range} bedeutet, dass der
Kanal nach Trendabzug im $s=-1$-Regime ruhiger bleibt und damit als
Regularisierungsmarker plausibler wird.
\end{itemize}

\subsection{Vergleich mit der Nullstellenseite}

Dieselben fuenf Observablen wurden zusaetzlich in
$10\,000$er-Fenstern gegen residuale Nullstellenkanaele verglichen:
\begin{itemize}
\item $\Delta_{\mathrm{res}}$,
\item $\Lambda_{\mathrm{res}}$,
\item $\Theta_{Z,\mathrm{res}}$.
\end{itemize}

Die zugehoerige Zusammenfassung liegt im Fenstervergleich zur
Nullstellenseite. Der zentrale Sammelindikator ist
\texttt{score\_abs\_sum}; je groesser dieser Wert ist, desto staerker ist
die kombinierte Kopplung einer Observable an die Nullstellenseite.

Damit liegen erstmals zwei Perspektiven zugleich in derselben
Auswertung vor:
\begin{itemize}
\item Regularisierungsruhe,
\item Nullstellenkopplung.
\end{itemize}

\subsection{Inhaltliches Zwischenurteil}

Der gegenwaertige numerische Stand laesst sich in vier Punkten
zusammenfassen:
\begin{enumerate}
\item Die Euler-Mascheroni-Konstante $\gamma$ gehoert tatsaechlich ins
Zentrum des $s=1$-Regimes. Die Darstellung
\[
\log\log X+\gamma
\]
erweist sich numerisch als sinnvolle Referenzachse.

\item Die additive und die multiplikative Seite trennen sich weiter
klar. $K_{\mathrm{add}}$ bleibt die natuerlichste additive Groesse,
waehrend $G_{P,235}$ weiterhin den staerksten multiplikativen
Haupttraeger bildet.

\item Die Brueckengroessen $\mathrm{Bridge}_{P,M}$ und
$\mathrm{Bridge}_{P,W}$ sind nicht nur dekorative Differenzen, sondern
plausible Kandidaten fuer eine vermittelnde Regularisierung zwischen dem
$\gamma$-zentrierten $s=1$-Bild und dem $s=-1$-regulierten additiven
Bild.

\item Der Abstand $d_e$ zur naechsten E-Primstelle ist inzwischen als
mitlaufende numerische Groesse fest eingebaut und kann als endlicher
Achsenabstieg des e-Anteils gegen den ABC-Kern gelesen werden.
\end{enumerate}

\subsection{Verdichtete Gesamtformulierung}

Die gegenwaertige Gesamtlage laesst sich wie folgt verdichten:

Die Euler-Mascheroni-Konstante $\gamma$ organisiert das harmonische
beziehungsweise polnahe $s=1$-Regime der Vierlingskanaele. Die additive
kleinsche Seite traegt das $s=-1$-Regime, waehrend die glatte
Produktschale die multiplikative Hauptstruktur traegt. Die
Brueckengroessen zwischen glattem und kleinschem Anteil sind die
natuerlichen Kandidaten fuer eine vermittelnde Regularisierung. Der
Abstand zur naechsten E-Primstelle fungiert dabei als endlicher
Achsenabstieg des e-Anteils gegen den ABC-Kern.

\subsection{Naechste Synthese}

Als naechster Schritt bietet sich eine kurze Ergebnis-Section oder ein
kompaktes Fazit an, das insbesondere folgende Fragen fokussiert:
\begin{itemize}
\item Welche Groesse ist derzeit der beste additive Marker?
\item Welche Groesse ist derzeit der beste multiplikative Marker?
\item Welche Bruecke vermittelt derzeit am besten?
\item Wo sitzen $\gamma$, $-1/12$ und der abc-Anschluss innerhalb des
gemeinsamen Bildes?
\end{itemize}

In dieser Form ist das Material als eigenstaendige, aber gut
kombinierbare LaTeX-Sektion vorbereitet.
 oder % Kombinierbarer Einschub fuer ein groesseres Hauptdokument.
% Vorgesehen fuer \input{Abschluss.tex} oder \include{Abschluss}.

\section{Numerische Unterlegung von Section 7}
\label{sec:section7_numerische_unterlegung}

Die bisherige inhaltliche Struktur von Section~7 laesst sich nun auch
numerisch unterlegen. Ziel ist es, die wesentlichen Marker nicht nur
qualitativ zu benennen, sondern ihre Stabilitaet in zwei verschiedenen
Regimen sowie ihre Kopplung an die Nullstellenseite sichtbar zu machen.

\subsection{Erzeugte Dateien}

Die numerische Auswertung fuehrt auf vier zentrale Ausgabedateien:
\begin{itemize}
\item \texttt{Section 7 -- Bericht},
\item \texttt{Alle Extrapolationswerte},
\item \texttt{Markeruebersicht fuer $s=1$ und $s=-1$},
\item \texttt{Fenstervergleich zur Nullstellenseite}.
\end{itemize}

Diese Dateien bilden zusammen die numerische Basis fuer die folgende
Interpretation.

\subsection{Untersuchte Observablen}

Betrachtet wurden fuenf zentrale Observablen:
\begin{itemize}
\item $K_{\mathrm{add}}$,
\item $G_{P,235}$,
\item $\mathrm{Bridge}_{P,M}$,
\item $\mathrm{Bridge}_{P,W}$,
\item $d_e$ als Abstand zur naechsten E-Primstelle.
\end{itemize}

Fuer jede dieser Groessen wurden zwei komplementaere Regime parallel
untersucht.

\subsection{Das $s=1$-Regime}

Im ersten Regime wurden die Partialsummen
\[
\sum_{p\le X}\frac{F(p)}{p}
\]
gegen
\[
\log\log X+\gamma
\]
gefittet. Die Euler-Mascheroni-Konstante
\[
\gamma \approx 0.5772156649
\]
fungiert dabei nicht nur als kosmetische Verschiebung, sondern als
explizites Zentrum des harmonischen beziehungsweise polnahen
$s=1$-Verhaltens.

\subsection{Das $s=-1$-Regime}

Im zweiten Regime wurden die rohen Summen
\[
\sum_{p\le X} F(p)\,p
\]
mit einem Trendmodell der Form
\[
a\,X\log X+b\,X+c\,\log X+d
\]
approximiert. Anschliessend wurden die Residuen untersucht. Diese
Regularisierung ist zwar noch keine analytische Fortsetzung im strengen
Sinn, sie ist numerisch aber ausreichend stabil, um zwischen ruhigen und
unruhigen Kanaelen zu unterscheiden.

\subsection{Lesart der Markeruebersicht}

Die zentrale Tabelle ist die Markeruebersicht fuer $s=1$ und $s=-1$.
Sie enthaelt fuer jede Observable insbesondere:
\begin{itemize}
\item die Steigung im $\gamma$-zentrierten $s=1$-Fit,
\item den letzten Residualwert,
\item die Gesamtbreite der Residuen,
\item die Koeffizienten des Trendmodells im $s=-1$-Regime.
\end{itemize}

Praktisch ist die Interpretation einfach:
\begin{itemize}
\item Eine kleine \texttt{s1\_residual\_range} bedeutet, dass sich der
Kanal gut an das Bild $\log\log X+\gamma$ anschliesst.
\item Eine kleine \texttt{sminus1\_residual\_range} bedeutet, dass der
Kanal nach Trendabzug im $s=-1$-Regime ruhiger bleibt und damit als
Regularisierungsmarker plausibler wird.
\end{itemize}

\subsection{Vergleich mit der Nullstellenseite}

Dieselben fuenf Observablen wurden zusaetzlich in
$10\,000$er-Fenstern gegen residuale Nullstellenkanaele verglichen:
\begin{itemize}
\item $\Delta_{\mathrm{res}}$,
\item $\Lambda_{\mathrm{res}}$,
\item $\Theta_{Z,\mathrm{res}}$.
\end{itemize}

Die zugehoerige Zusammenfassung liegt im Fenstervergleich zur
Nullstellenseite. Der zentrale Sammelindikator ist
\texttt{score\_abs\_sum}; je groesser dieser Wert ist, desto staerker ist
die kombinierte Kopplung einer Observable an die Nullstellenseite.

Damit liegen erstmals zwei Perspektiven zugleich in derselben
Auswertung vor:
\begin{itemize}
\item Regularisierungsruhe,
\item Nullstellenkopplung.
\end{itemize}

\subsection{Inhaltliches Zwischenurteil}

Der gegenwaertige numerische Stand laesst sich in vier Punkten
zusammenfassen:
\begin{enumerate}
\item Die Euler-Mascheroni-Konstante $\gamma$ gehoert tatsaechlich ins
Zentrum des $s=1$-Regimes. Die Darstellung
\[
\log\log X+\gamma
\]
erweist sich numerisch als sinnvolle Referenzachse.

\item Die additive und die multiplikative Seite trennen sich weiter
klar. $K_{\mathrm{add}}$ bleibt die natuerlichste additive Groesse,
waehrend $G_{P,235}$ weiterhin den staerksten multiplikativen
Haupttraeger bildet.

\item Die Brueckengroessen $\mathrm{Bridge}_{P,M}$ und
$\mathrm{Bridge}_{P,W}$ sind nicht nur dekorative Differenzen, sondern
plausible Kandidaten fuer eine vermittelnde Regularisierung zwischen dem
$\gamma$-zentrierten $s=1$-Bild und dem $s=-1$-regulierten additiven
Bild.

\item Der Abstand $d_e$ zur naechsten E-Primstelle ist inzwischen als
mitlaufende numerische Groesse fest eingebaut und kann als endlicher
Achsenabstieg des e-Anteils gegen den ABC-Kern gelesen werden.
\end{enumerate}

\subsection{Verdichtete Gesamtformulierung}

Die gegenwaertige Gesamtlage laesst sich wie folgt verdichten:

Die Euler-Mascheroni-Konstante $\gamma$ organisiert das harmonische
beziehungsweise polnahe $s=1$-Regime der Vierlingskanaele. Die additive
kleinsche Seite traegt das $s=-1$-Regime, waehrend die glatte
Produktschale die multiplikative Hauptstruktur traegt. Die
Brueckengroessen zwischen glattem und kleinschem Anteil sind die
natuerlichen Kandidaten fuer eine vermittelnde Regularisierung. Der
Abstand zur naechsten E-Primstelle fungiert dabei als endlicher
Achsenabstieg des e-Anteils gegen den ABC-Kern.

\subsection{Naechste Synthese}

Als naechster Schritt bietet sich eine kurze Ergebnis-Section oder ein
kompaktes Fazit an, das insbesondere folgende Fragen fokussiert:
\begin{itemize}
\item Welche Groesse ist derzeit der beste additive Marker?
\item Welche Groesse ist derzeit der beste multiplikative Marker?
\item Welche Bruecke vermittelt derzeit am besten?
\item Wo sitzen $\gamma$, $-1/12$ und der abc-Anschluss innerhalb des
gemeinsamen Bildes?
\end{itemize}

In dieser Form ist das Material als eigenstaendige, aber gut
kombinierbare LaTeX-Sektion vorbereitet.
.

\section{Numerische Unterlegung von Section 7}
\label{sec:section7_numerische_unterlegung}

Die bisherige inhaltliche Struktur von Section~7 laesst sich nun auch
numerisch unterlegen. Ziel ist es, die wesentlichen Marker nicht nur
qualitativ zu benennen, sondern ihre Stabilitaet in zwei verschiedenen
Regimen sowie ihre Kopplung an die Nullstellenseite sichtbar zu machen.

\subsection{Erzeugte Dateien}

Die numerische Auswertung fuehrt auf vier zentrale Ausgabedateien:
\begin{itemize}
\item \texttt{Section 7 -- Bericht},
\item \texttt{Alle Extrapolationswerte},
\item \texttt{Markeruebersicht fuer $s=1$ und $s=-1$},
\item \texttt{Fenstervergleich zur Nullstellenseite}.
\end{itemize}

Diese Dateien bilden zusammen die numerische Basis fuer die folgende
Interpretation.

\subsection{Untersuchte Observablen}

Betrachtet wurden fuenf zentrale Observablen:
\begin{itemize}
\item $K_{\mathrm{add}}$,
\item $G_{P,235}$,
\item $\mathrm{Bridge}_{P,M}$,
\item $\mathrm{Bridge}_{P,W}$,
\item $d_e$ als Abstand zur naechsten E-Primstelle.
\end{itemize}

Fuer jede dieser Groessen wurden zwei komplementaere Regime parallel
untersucht.

\subsection{Das $s=1$-Regime}

Im ersten Regime wurden die Partialsummen
\[
\sum_{p\le X}\frac{F(p)}{p}
\]
gegen
\[
\log\log X+\gamma
\]
gefittet. Die Euler-Mascheroni-Konstante
\[
\gamma \approx 0.5772156649
\]
fungiert dabei nicht nur als kosmetische Verschiebung, sondern als
explizites Zentrum des harmonischen beziehungsweise polnahen
$s=1$-Verhaltens.

\subsection{Das $s=-1$-Regime}

Im zweiten Regime wurden die rohen Summen
\[
\sum_{p\le X} F(p)\,p
\]
mit einem Trendmodell der Form
\[
a\,X\log X+b\,X+c\,\log X+d
\]
approximiert. Anschliessend wurden die Residuen untersucht. Diese
Regularisierung ist zwar noch keine analytische Fortsetzung im strengen
Sinn, sie ist numerisch aber ausreichend stabil, um zwischen ruhigen und
unruhigen Kanaelen zu unterscheiden.

\subsection{Lesart der Markeruebersicht}

Die zentrale Tabelle ist die Markeruebersicht fuer $s=1$ und $s=-1$.
Sie enthaelt fuer jede Observable insbesondere:
\begin{itemize}
\item die Steigung im $\gamma$-zentrierten $s=1$-Fit,
\item den letzten Residualwert,
\item die Gesamtbreite der Residuen,
\item die Koeffizienten des Trendmodells im $s=-1$-Regime.
\end{itemize}

Praktisch ist die Interpretation einfach:
\begin{itemize}
\item Eine kleine \texttt{s1\_residual\_range} bedeutet, dass sich der
Kanal gut an das Bild $\log\log X+\gamma$ anschliesst.
\item Eine kleine \texttt{sminus1\_residual\_range} bedeutet, dass der
Kanal nach Trendabzug im $s=-1$-Regime ruhiger bleibt und damit als
Regularisierungsmarker plausibler wird.
\end{itemize}

\subsection{Vergleich mit der Nullstellenseite}

Dieselben fuenf Observablen wurden zusaetzlich in
$10\,000$er-Fenstern gegen residuale Nullstellenkanaele verglichen:
\begin{itemize}
\item $\Delta_{\mathrm{res}}$,
\item $\Lambda_{\mathrm{res}}$,
\item $\Theta_{Z,\mathrm{res}}$.
\end{itemize}

Die zugehoerige Zusammenfassung liegt im Fenstervergleich zur
Nullstellenseite. Der zentrale Sammelindikator ist
\texttt{score\_abs\_sum}; je groesser dieser Wert ist, desto staerker ist
die kombinierte Kopplung einer Observable an die Nullstellenseite.

Damit liegen erstmals zwei Perspektiven zugleich in derselben
Auswertung vor:
\begin{itemize}
\item Regularisierungsruhe,
\item Nullstellenkopplung.
\end{itemize}

\subsection{Inhaltliches Zwischenurteil}

Der gegenwaertige numerische Stand laesst sich in vier Punkten
zusammenfassen:
\begin{enumerate}
\item Die Euler-Mascheroni-Konstante $\gamma$ gehoert tatsaechlich ins
Zentrum des $s=1$-Regimes. Die Darstellung
\[
\log\log X+\gamma
\]
erweist sich numerisch als sinnvolle Referenzachse.

\item Die additive und die multiplikative Seite trennen sich weiter
klar. $K_{\mathrm{add}}$ bleibt die natuerlichste additive Groesse,
waehrend $G_{P,235}$ weiterhin den staerksten multiplikativen
Haupttraeger bildet.

\item Die Brueckengroessen $\mathrm{Bridge}_{P,M}$ und
$\mathrm{Bridge}_{P,W}$ sind nicht nur dekorative Differenzen, sondern
plausible Kandidaten fuer eine vermittelnde Regularisierung zwischen dem
$\gamma$-zentrierten $s=1$-Bild und dem $s=-1$-regulierten additiven
Bild.

\item Der Abstand $d_e$ zur naechsten E-Primstelle ist inzwischen als
mitlaufende numerische Groesse fest eingebaut und kann als endlicher
Achsenabstieg des e-Anteils gegen den ABC-Kern gelesen werden.
\end{enumerate}

\subsection{Verdichtete Gesamtformulierung}

Die gegenwaertige Gesamtlage laesst sich wie folgt verdichten:

Die Euler-Mascheroni-Konstante $\gamma$ organisiert das harmonische
beziehungsweise polnahe $s=1$-Regime der Vierlingskanaele. Die additive
kleinsche Seite traegt das $s=-1$-Regime, waehrend die glatte
Produktschale die multiplikative Hauptstruktur traegt. Die
Brueckengroessen zwischen glattem und kleinschem Anteil sind die
natuerlichen Kandidaten fuer eine vermittelnde Regularisierung. Der
Abstand zur naechsten E-Primstelle fungiert dabei als endlicher
Achsenabstieg des e-Anteils gegen den ABC-Kern.

\subsection{Naechste Synthese}

Als naechster Schritt bietet sich eine kurze Ergebnis-Section oder ein
kompaktes Fazit an, das insbesondere folgende Fragen fokussiert:
\begin{itemize}
\item Welche Groesse ist derzeit der beste additive Marker?
\item Welche Groesse ist derzeit der beste multiplikative Marker?
\item Welche Bruecke vermittelt derzeit am besten?
\item Wo sitzen $\gamma$, $-1/12$ und der abc-Anschluss innerhalb des
gemeinsamen Bildes?
\end{itemize}

In dieser Form ist das Material als eigenstaendige, aber gut
kombinierbare LaTeX-Sektion vorbereitet.
 oder % Kombinierbarer Einschub fuer ein groesseres Hauptdokument.
% Vorgesehen fuer % Kombinierbarer Einschub fuer ein groesseres Hauptdokument.
% Vorgesehen fuer \input{Abschluss.tex} oder \include{Abschluss}.

\section{Numerische Unterlegung von Section 7}
\label{sec:section7_numerische_unterlegung}

Die bisherige inhaltliche Struktur von Section~7 laesst sich nun auch
numerisch unterlegen. Ziel ist es, die wesentlichen Marker nicht nur
qualitativ zu benennen, sondern ihre Stabilitaet in zwei verschiedenen
Regimen sowie ihre Kopplung an die Nullstellenseite sichtbar zu machen.

\subsection{Erzeugte Dateien}

Die numerische Auswertung fuehrt auf vier zentrale Ausgabedateien:
\begin{itemize}
\item \texttt{Section 7 -- Bericht},
\item \texttt{Alle Extrapolationswerte},
\item \texttt{Markeruebersicht fuer $s=1$ und $s=-1$},
\item \texttt{Fenstervergleich zur Nullstellenseite}.
\end{itemize}

Diese Dateien bilden zusammen die numerische Basis fuer die folgende
Interpretation.

\subsection{Untersuchte Observablen}

Betrachtet wurden fuenf zentrale Observablen:
\begin{itemize}
\item $K_{\mathrm{add}}$,
\item $G_{P,235}$,
\item $\mathrm{Bridge}_{P,M}$,
\item $\mathrm{Bridge}_{P,W}$,
\item $d_e$ als Abstand zur naechsten E-Primstelle.
\end{itemize}

Fuer jede dieser Groessen wurden zwei komplementaere Regime parallel
untersucht.

\subsection{Das $s=1$-Regime}

Im ersten Regime wurden die Partialsummen
\[
\sum_{p\le X}\frac{F(p)}{p}
\]
gegen
\[
\log\log X+\gamma
\]
gefittet. Die Euler-Mascheroni-Konstante
\[
\gamma \approx 0.5772156649
\]
fungiert dabei nicht nur als kosmetische Verschiebung, sondern als
explizites Zentrum des harmonischen beziehungsweise polnahen
$s=1$-Verhaltens.

\subsection{Das $s=-1$-Regime}

Im zweiten Regime wurden die rohen Summen
\[
\sum_{p\le X} F(p)\,p
\]
mit einem Trendmodell der Form
\[
a\,X\log X+b\,X+c\,\log X+d
\]
approximiert. Anschliessend wurden die Residuen untersucht. Diese
Regularisierung ist zwar noch keine analytische Fortsetzung im strengen
Sinn, sie ist numerisch aber ausreichend stabil, um zwischen ruhigen und
unruhigen Kanaelen zu unterscheiden.

\subsection{Lesart der Markeruebersicht}

Die zentrale Tabelle ist die Markeruebersicht fuer $s=1$ und $s=-1$.
Sie enthaelt fuer jede Observable insbesondere:
\begin{itemize}
\item die Steigung im $\gamma$-zentrierten $s=1$-Fit,
\item den letzten Residualwert,
\item die Gesamtbreite der Residuen,
\item die Koeffizienten des Trendmodells im $s=-1$-Regime.
\end{itemize}

Praktisch ist die Interpretation einfach:
\begin{itemize}
\item Eine kleine \texttt{s1\_residual\_range} bedeutet, dass sich der
Kanal gut an das Bild $\log\log X+\gamma$ anschliesst.
\item Eine kleine \texttt{sminus1\_residual\_range} bedeutet, dass der
Kanal nach Trendabzug im $s=-1$-Regime ruhiger bleibt und damit als
Regularisierungsmarker plausibler wird.
\end{itemize}

\subsection{Vergleich mit der Nullstellenseite}

Dieselben fuenf Observablen wurden zusaetzlich in
$10\,000$er-Fenstern gegen residuale Nullstellenkanaele verglichen:
\begin{itemize}
\item $\Delta_{\mathrm{res}}$,
\item $\Lambda_{\mathrm{res}}$,
\item $\Theta_{Z,\mathrm{res}}$.
\end{itemize}

Die zugehoerige Zusammenfassung liegt im Fenstervergleich zur
Nullstellenseite. Der zentrale Sammelindikator ist
\texttt{score\_abs\_sum}; je groesser dieser Wert ist, desto staerker ist
die kombinierte Kopplung einer Observable an die Nullstellenseite.

Damit liegen erstmals zwei Perspektiven zugleich in derselben
Auswertung vor:
\begin{itemize}
\item Regularisierungsruhe,
\item Nullstellenkopplung.
\end{itemize}

\subsection{Inhaltliches Zwischenurteil}

Der gegenwaertige numerische Stand laesst sich in vier Punkten
zusammenfassen:
\begin{enumerate}
\item Die Euler-Mascheroni-Konstante $\gamma$ gehoert tatsaechlich ins
Zentrum des $s=1$-Regimes. Die Darstellung
\[
\log\log X+\gamma
\]
erweist sich numerisch als sinnvolle Referenzachse.

\item Die additive und die multiplikative Seite trennen sich weiter
klar. $K_{\mathrm{add}}$ bleibt die natuerlichste additive Groesse,
waehrend $G_{P,235}$ weiterhin den staerksten multiplikativen
Haupttraeger bildet.

\item Die Brueckengroessen $\mathrm{Bridge}_{P,M}$ und
$\mathrm{Bridge}_{P,W}$ sind nicht nur dekorative Differenzen, sondern
plausible Kandidaten fuer eine vermittelnde Regularisierung zwischen dem
$\gamma$-zentrierten $s=1$-Bild und dem $s=-1$-regulierten additiven
Bild.

\item Der Abstand $d_e$ zur naechsten E-Primstelle ist inzwischen als
mitlaufende numerische Groesse fest eingebaut und kann als endlicher
Achsenabstieg des e-Anteils gegen den ABC-Kern gelesen werden.
\end{enumerate}

\subsection{Verdichtete Gesamtformulierung}

Die gegenwaertige Gesamtlage laesst sich wie folgt verdichten:

Die Euler-Mascheroni-Konstante $\gamma$ organisiert das harmonische
beziehungsweise polnahe $s=1$-Regime der Vierlingskanaele. Die additive
kleinsche Seite traegt das $s=-1$-Regime, waehrend die glatte
Produktschale die multiplikative Hauptstruktur traegt. Die
Brueckengroessen zwischen glattem und kleinschem Anteil sind die
natuerlichen Kandidaten fuer eine vermittelnde Regularisierung. Der
Abstand zur naechsten E-Primstelle fungiert dabei als endlicher
Achsenabstieg des e-Anteils gegen den ABC-Kern.

\subsection{Naechste Synthese}

Als naechster Schritt bietet sich eine kurze Ergebnis-Section oder ein
kompaktes Fazit an, das insbesondere folgende Fragen fokussiert:
\begin{itemize}
\item Welche Groesse ist derzeit der beste additive Marker?
\item Welche Groesse ist derzeit der beste multiplikative Marker?
\item Welche Bruecke vermittelt derzeit am besten?
\item Wo sitzen $\gamma$, $-1/12$ und der abc-Anschluss innerhalb des
gemeinsamen Bildes?
\end{itemize}

In dieser Form ist das Material als eigenstaendige, aber gut
kombinierbare LaTeX-Sektion vorbereitet.
 oder % Kombinierbarer Einschub fuer ein groesseres Hauptdokument.
% Vorgesehen fuer \input{Abschluss.tex} oder \include{Abschluss}.

\section{Numerische Unterlegung von Section 7}
\label{sec:section7_numerische_unterlegung}

Die bisherige inhaltliche Struktur von Section~7 laesst sich nun auch
numerisch unterlegen. Ziel ist es, die wesentlichen Marker nicht nur
qualitativ zu benennen, sondern ihre Stabilitaet in zwei verschiedenen
Regimen sowie ihre Kopplung an die Nullstellenseite sichtbar zu machen.

\subsection{Erzeugte Dateien}

Die numerische Auswertung fuehrt auf vier zentrale Ausgabedateien:
\begin{itemize}
\item \texttt{Section 7 -- Bericht},
\item \texttt{Alle Extrapolationswerte},
\item \texttt{Markeruebersicht fuer $s=1$ und $s=-1$},
\item \texttt{Fenstervergleich zur Nullstellenseite}.
\end{itemize}

Diese Dateien bilden zusammen die numerische Basis fuer die folgende
Interpretation.

\subsection{Untersuchte Observablen}

Betrachtet wurden fuenf zentrale Observablen:
\begin{itemize}
\item $K_{\mathrm{add}}$,
\item $G_{P,235}$,
\item $\mathrm{Bridge}_{P,M}$,
\item $\mathrm{Bridge}_{P,W}$,
\item $d_e$ als Abstand zur naechsten E-Primstelle.
\end{itemize}

Fuer jede dieser Groessen wurden zwei komplementaere Regime parallel
untersucht.

\subsection{Das $s=1$-Regime}

Im ersten Regime wurden die Partialsummen
\[
\sum_{p\le X}\frac{F(p)}{p}
\]
gegen
\[
\log\log X+\gamma
\]
gefittet. Die Euler-Mascheroni-Konstante
\[
\gamma \approx 0.5772156649
\]
fungiert dabei nicht nur als kosmetische Verschiebung, sondern als
explizites Zentrum des harmonischen beziehungsweise polnahen
$s=1$-Verhaltens.

\subsection{Das $s=-1$-Regime}

Im zweiten Regime wurden die rohen Summen
\[
\sum_{p\le X} F(p)\,p
\]
mit einem Trendmodell der Form
\[
a\,X\log X+b\,X+c\,\log X+d
\]
approximiert. Anschliessend wurden die Residuen untersucht. Diese
Regularisierung ist zwar noch keine analytische Fortsetzung im strengen
Sinn, sie ist numerisch aber ausreichend stabil, um zwischen ruhigen und
unruhigen Kanaelen zu unterscheiden.

\subsection{Lesart der Markeruebersicht}

Die zentrale Tabelle ist die Markeruebersicht fuer $s=1$ und $s=-1$.
Sie enthaelt fuer jede Observable insbesondere:
\begin{itemize}
\item die Steigung im $\gamma$-zentrierten $s=1$-Fit,
\item den letzten Residualwert,
\item die Gesamtbreite der Residuen,
\item die Koeffizienten des Trendmodells im $s=-1$-Regime.
\end{itemize}

Praktisch ist die Interpretation einfach:
\begin{itemize}
\item Eine kleine \texttt{s1\_residual\_range} bedeutet, dass sich der
Kanal gut an das Bild $\log\log X+\gamma$ anschliesst.
\item Eine kleine \texttt{sminus1\_residual\_range} bedeutet, dass der
Kanal nach Trendabzug im $s=-1$-Regime ruhiger bleibt und damit als
Regularisierungsmarker plausibler wird.
\end{itemize}

\subsection{Vergleich mit der Nullstellenseite}

Dieselben fuenf Observablen wurden zusaetzlich in
$10\,000$er-Fenstern gegen residuale Nullstellenkanaele verglichen:
\begin{itemize}
\item $\Delta_{\mathrm{res}}$,
\item $\Lambda_{\mathrm{res}}$,
\item $\Theta_{Z,\mathrm{res}}$.
\end{itemize}

Die zugehoerige Zusammenfassung liegt im Fenstervergleich zur
Nullstellenseite. Der zentrale Sammelindikator ist
\texttt{score\_abs\_sum}; je groesser dieser Wert ist, desto staerker ist
die kombinierte Kopplung einer Observable an die Nullstellenseite.

Damit liegen erstmals zwei Perspektiven zugleich in derselben
Auswertung vor:
\begin{itemize}
\item Regularisierungsruhe,
\item Nullstellenkopplung.
\end{itemize}

\subsection{Inhaltliches Zwischenurteil}

Der gegenwaertige numerische Stand laesst sich in vier Punkten
zusammenfassen:
\begin{enumerate}
\item Die Euler-Mascheroni-Konstante $\gamma$ gehoert tatsaechlich ins
Zentrum des $s=1$-Regimes. Die Darstellung
\[
\log\log X+\gamma
\]
erweist sich numerisch als sinnvolle Referenzachse.

\item Die additive und die multiplikative Seite trennen sich weiter
klar. $K_{\mathrm{add}}$ bleibt die natuerlichste additive Groesse,
waehrend $G_{P,235}$ weiterhin den staerksten multiplikativen
Haupttraeger bildet.

\item Die Brueckengroessen $\mathrm{Bridge}_{P,M}$ und
$\mathrm{Bridge}_{P,W}$ sind nicht nur dekorative Differenzen, sondern
plausible Kandidaten fuer eine vermittelnde Regularisierung zwischen dem
$\gamma$-zentrierten $s=1$-Bild und dem $s=-1$-regulierten additiven
Bild.

\item Der Abstand $d_e$ zur naechsten E-Primstelle ist inzwischen als
mitlaufende numerische Groesse fest eingebaut und kann als endlicher
Achsenabstieg des e-Anteils gegen den ABC-Kern gelesen werden.
\end{enumerate}

\subsection{Verdichtete Gesamtformulierung}

Die gegenwaertige Gesamtlage laesst sich wie folgt verdichten:

Die Euler-Mascheroni-Konstante $\gamma$ organisiert das harmonische
beziehungsweise polnahe $s=1$-Regime der Vierlingskanaele. Die additive
kleinsche Seite traegt das $s=-1$-Regime, waehrend die glatte
Produktschale die multiplikative Hauptstruktur traegt. Die
Brueckengroessen zwischen glattem und kleinschem Anteil sind die
natuerlichen Kandidaten fuer eine vermittelnde Regularisierung. Der
Abstand zur naechsten E-Primstelle fungiert dabei als endlicher
Achsenabstieg des e-Anteils gegen den ABC-Kern.

\subsection{Naechste Synthese}

Als naechster Schritt bietet sich eine kurze Ergebnis-Section oder ein
kompaktes Fazit an, das insbesondere folgende Fragen fokussiert:
\begin{itemize}
\item Welche Groesse ist derzeit der beste additive Marker?
\item Welche Groesse ist derzeit der beste multiplikative Marker?
\item Welche Bruecke vermittelt derzeit am besten?
\item Wo sitzen $\gamma$, $-1/12$ und der abc-Anschluss innerhalb des
gemeinsamen Bildes?
\end{itemize}

In dieser Form ist das Material als eigenstaendige, aber gut
kombinierbare LaTeX-Sektion vorbereitet.
.

\section{Numerische Unterlegung von Section 7}
\label{sec:section7_numerische_unterlegung}

Die bisherige inhaltliche Struktur von Section~7 laesst sich nun auch
numerisch unterlegen. Ziel ist es, die wesentlichen Marker nicht nur
qualitativ zu benennen, sondern ihre Stabilitaet in zwei verschiedenen
Regimen sowie ihre Kopplung an die Nullstellenseite sichtbar zu machen.

\subsection{Erzeugte Dateien}

Die numerische Auswertung fuehrt auf vier zentrale Ausgabedateien:
\begin{itemize}
\item \texttt{Section 7 -- Bericht},
\item \texttt{Alle Extrapolationswerte},
\item \texttt{Markeruebersicht fuer $s=1$ und $s=-1$},
\item \texttt{Fenstervergleich zur Nullstellenseite}.
\end{itemize}

Diese Dateien bilden zusammen die numerische Basis fuer die folgende
Interpretation.

\subsection{Untersuchte Observablen}

Betrachtet wurden fuenf zentrale Observablen:
\begin{itemize}
\item $K_{\mathrm{add}}$,
\item $G_{P,235}$,
\item $\mathrm{Bridge}_{P,M}$,
\item $\mathrm{Bridge}_{P,W}$,
\item $d_e$ als Abstand zur naechsten E-Primstelle.
\end{itemize}

Fuer jede dieser Groessen wurden zwei komplementaere Regime parallel
untersucht.

\subsection{Das $s=1$-Regime}

Im ersten Regime wurden die Partialsummen
\[
\sum_{p\le X}\frac{F(p)}{p}
\]
gegen
\[
\log\log X+\gamma
\]
gefittet. Die Euler-Mascheroni-Konstante
\[
\gamma \approx 0.5772156649
\]
fungiert dabei nicht nur als kosmetische Verschiebung, sondern als
explizites Zentrum des harmonischen beziehungsweise polnahen
$s=1$-Verhaltens.

\subsection{Das $s=-1$-Regime}

Im zweiten Regime wurden die rohen Summen
\[
\sum_{p\le X} F(p)\,p
\]
mit einem Trendmodell der Form
\[
a\,X\log X+b\,X+c\,\log X+d
\]
approximiert. Anschliessend wurden die Residuen untersucht. Diese
Regularisierung ist zwar noch keine analytische Fortsetzung im strengen
Sinn, sie ist numerisch aber ausreichend stabil, um zwischen ruhigen und
unruhigen Kanaelen zu unterscheiden.

\subsection{Lesart der Markeruebersicht}

Die zentrale Tabelle ist die Markeruebersicht fuer $s=1$ und $s=-1$.
Sie enthaelt fuer jede Observable insbesondere:
\begin{itemize}
\item die Steigung im $\gamma$-zentrierten $s=1$-Fit,
\item den letzten Residualwert,
\item die Gesamtbreite der Residuen,
\item die Koeffizienten des Trendmodells im $s=-1$-Regime.
\end{itemize}

Praktisch ist die Interpretation einfach:
\begin{itemize}
\item Eine kleine \texttt{s1\_residual\_range} bedeutet, dass sich der
Kanal gut an das Bild $\log\log X+\gamma$ anschliesst.
\item Eine kleine \texttt{sminus1\_residual\_range} bedeutet, dass der
Kanal nach Trendabzug im $s=-1$-Regime ruhiger bleibt und damit als
Regularisierungsmarker plausibler wird.
\end{itemize}

\subsection{Vergleich mit der Nullstellenseite}

Dieselben fuenf Observablen wurden zusaetzlich in
$10\,000$er-Fenstern gegen residuale Nullstellenkanaele verglichen:
\begin{itemize}
\item $\Delta_{\mathrm{res}}$,
\item $\Lambda_{\mathrm{res}}$,
\item $\Theta_{Z,\mathrm{res}}$.
\end{itemize}

Die zugehoerige Zusammenfassung liegt im Fenstervergleich zur
Nullstellenseite. Der zentrale Sammelindikator ist
\texttt{score\_abs\_sum}; je groesser dieser Wert ist, desto staerker ist
die kombinierte Kopplung einer Observable an die Nullstellenseite.

Damit liegen erstmals zwei Perspektiven zugleich in derselben
Auswertung vor:
\begin{itemize}
\item Regularisierungsruhe,
\item Nullstellenkopplung.
\end{itemize}

\subsection{Inhaltliches Zwischenurteil}

Der gegenwaertige numerische Stand laesst sich in vier Punkten
zusammenfassen:
\begin{enumerate}
\item Die Euler-Mascheroni-Konstante $\gamma$ gehoert tatsaechlich ins
Zentrum des $s=1$-Regimes. Die Darstellung
\[
\log\log X+\gamma
\]
erweist sich numerisch als sinnvolle Referenzachse.

\item Die additive und die multiplikative Seite trennen sich weiter
klar. $K_{\mathrm{add}}$ bleibt die natuerlichste additive Groesse,
waehrend $G_{P,235}$ weiterhin den staerksten multiplikativen
Haupttraeger bildet.

\item Die Brueckengroessen $\mathrm{Bridge}_{P,M}$ und
$\mathrm{Bridge}_{P,W}$ sind nicht nur dekorative Differenzen, sondern
plausible Kandidaten fuer eine vermittelnde Regularisierung zwischen dem
$\gamma$-zentrierten $s=1$-Bild und dem $s=-1$-regulierten additiven
Bild.

\item Der Abstand $d_e$ zur naechsten E-Primstelle ist inzwischen als
mitlaufende numerische Groesse fest eingebaut und kann als endlicher
Achsenabstieg des e-Anteils gegen den ABC-Kern gelesen werden.
\end{enumerate}

\subsection{Verdichtete Gesamtformulierung}

Die gegenwaertige Gesamtlage laesst sich wie folgt verdichten:

Die Euler-Mascheroni-Konstante $\gamma$ organisiert das harmonische
beziehungsweise polnahe $s=1$-Regime der Vierlingskanaele. Die additive
kleinsche Seite traegt das $s=-1$-Regime, waehrend die glatte
Produktschale die multiplikative Hauptstruktur traegt. Die
Brueckengroessen zwischen glattem und kleinschem Anteil sind die
natuerlichen Kandidaten fuer eine vermittelnde Regularisierung. Der
Abstand zur naechsten E-Primstelle fungiert dabei als endlicher
Achsenabstieg des e-Anteils gegen den ABC-Kern.

\subsection{Naechste Synthese}

Als naechster Schritt bietet sich eine kurze Ergebnis-Section oder ein
kompaktes Fazit an, das insbesondere folgende Fragen fokussiert:
\begin{itemize}
\item Welche Groesse ist derzeit der beste additive Marker?
\item Welche Groesse ist derzeit der beste multiplikative Marker?
\item Welche Bruecke vermittelt derzeit am besten?
\item Wo sitzen $\gamma$, $-1/12$ und der abc-Anschluss innerhalb des
gemeinsamen Bildes?
\end{itemize}

In dieser Form ist das Material als eigenstaendige, aber gut
kombinierbare LaTeX-Sektion vorbereitet.
.

\section{Numerische Unterlegung von Section 7}
\label{sec:section7_numerische_unterlegung}

Die bisherige inhaltliche Struktur von Section~7 laesst sich nun auch
numerisch unterlegen. Ziel ist es, die wesentlichen Marker nicht nur
qualitativ zu benennen, sondern ihre Stabilitaet in zwei verschiedenen
Regimen sowie ihre Kopplung an die Nullstellenseite sichtbar zu machen.

\subsection{Erzeugte Dateien}

Die numerische Auswertung fuehrt auf vier zentrale Ausgabedateien:
\begin{itemize}
\item \texttt{Section 7 -- Bericht},
\item \texttt{Alle Extrapolationswerte},
\item \texttt{Markeruebersicht fuer $s=1$ und $s=-1$},
\item \texttt{Fenstervergleich zur Nullstellenseite}.
\end{itemize}

Diese Dateien bilden zusammen die numerische Basis fuer die folgende
Interpretation.

\subsection{Untersuchte Observablen}

Betrachtet wurden fuenf zentrale Observablen:
\begin{itemize}
\item $K_{\mathrm{add}}$,
\item $G_{P,235}$,
\item $\mathrm{Bridge}_{P,M}$,
\item $\mathrm{Bridge}_{P,W}$,
\item $d_e$ als Abstand zur naechsten E-Primstelle.
\end{itemize}

Fuer jede dieser Groessen wurden zwei komplementaere Regime parallel
untersucht.

\subsection{Das $s=1$-Regime}

Im ersten Regime wurden die Partialsummen
\[
\sum_{p\le X}\frac{F(p)}{p}
\]
gegen
\[
\log\log X+\gamma
\]
gefittet. Die Euler-Mascheroni-Konstante
\[
\gamma \approx 0.5772156649
\]
fungiert dabei nicht nur als kosmetische Verschiebung, sondern als
explizites Zentrum des harmonischen beziehungsweise polnahen
$s=1$-Verhaltens.

\subsection{Das $s=-1$-Regime}

Im zweiten Regime wurden die rohen Summen
\[
\sum_{p\le X} F(p)\,p
\]
mit einem Trendmodell der Form
\[
a\,X\log X+b\,X+c\,\log X+d
\]
approximiert. Anschliessend wurden die Residuen untersucht. Diese
Regularisierung ist zwar noch keine analytische Fortsetzung im strengen
Sinn, sie ist numerisch aber ausreichend stabil, um zwischen ruhigen und
unruhigen Kanaelen zu unterscheiden.

\subsection{Lesart der Markeruebersicht}

Die zentrale Tabelle ist die Markeruebersicht fuer $s=1$ und $s=-1$.
Sie enthaelt fuer jede Observable insbesondere:
\begin{itemize}
\item die Steigung im $\gamma$-zentrierten $s=1$-Fit,
\item den letzten Residualwert,
\item die Gesamtbreite der Residuen,
\item die Koeffizienten des Trendmodells im $s=-1$-Regime.
\end{itemize}

Praktisch ist die Interpretation einfach:
\begin{itemize}
\item Eine kleine \texttt{s1\_residual\_range} bedeutet, dass sich der
Kanal gut an das Bild $\log\log X+\gamma$ anschliesst.
\item Eine kleine \texttt{sminus1\_residual\_range} bedeutet, dass der
Kanal nach Trendabzug im $s=-1$-Regime ruhiger bleibt und damit als
Regularisierungsmarker plausibler wird.
\end{itemize}

\subsection{Vergleich mit der Nullstellenseite}

Dieselben fuenf Observablen wurden zusaetzlich in
$10\,000$er-Fenstern gegen residuale Nullstellenkanaele verglichen:
\begin{itemize}
\item $\Delta_{\mathrm{res}}$,
\item $\Lambda_{\mathrm{res}}$,
\item $\Theta_{Z,\mathrm{res}}$.
\end{itemize}

Die zugehoerige Zusammenfassung liegt im Fenstervergleich zur
Nullstellenseite. Der zentrale Sammelindikator ist
\texttt{score\_abs\_sum}; je groesser dieser Wert ist, desto staerker ist
die kombinierte Kopplung einer Observable an die Nullstellenseite.

Damit liegen erstmals zwei Perspektiven zugleich in derselben
Auswertung vor:
\begin{itemize}
\item Regularisierungsruhe,
\item Nullstellenkopplung.
\end{itemize}

\subsection{Inhaltliches Zwischenurteil}

Der gegenwaertige numerische Stand laesst sich in vier Punkten
zusammenfassen:
\begin{enumerate}
\item Die Euler-Mascheroni-Konstante $\gamma$ gehoert tatsaechlich ins
Zentrum des $s=1$-Regimes. Die Darstellung
\[
\log\log X+\gamma
\]
erweist sich numerisch als sinnvolle Referenzachse.

\item Die additive und die multiplikative Seite trennen sich weiter
klar. $K_{\mathrm{add}}$ bleibt die natuerlichste additive Groesse,
waehrend $G_{P,235}$ weiterhin den staerksten multiplikativen
Haupttraeger bildet.

\item Die Brueckengroessen $\mathrm{Bridge}_{P,M}$ und
$\mathrm{Bridge}_{P,W}$ sind nicht nur dekorative Differenzen, sondern
plausible Kandidaten fuer eine vermittelnde Regularisierung zwischen dem
$\gamma$-zentrierten $s=1$-Bild und dem $s=-1$-regulierten additiven
Bild.

\item Der Abstand $d_e$ zur naechsten E-Primstelle ist inzwischen als
mitlaufende numerische Groesse fest eingebaut und kann als endlicher
Achsenabstieg des e-Anteils gegen den ABC-Kern gelesen werden.
\end{enumerate}

\subsection{Verdichtete Gesamtformulierung}

Die gegenwaertige Gesamtlage laesst sich wie folgt verdichten:

Die Euler-Mascheroni-Konstante $\gamma$ organisiert das harmonische
beziehungsweise polnahe $s=1$-Regime der Vierlingskanaele. Die additive
kleinsche Seite traegt das $s=-1$-Regime, waehrend die glatte
Produktschale die multiplikative Hauptstruktur traegt. Die
Brueckengroessen zwischen glattem und kleinschem Anteil sind die
natuerlichen Kandidaten fuer eine vermittelnde Regularisierung. Der
Abstand zur naechsten E-Primstelle fungiert dabei als endlicher
Achsenabstieg des e-Anteils gegen den ABC-Kern.

\subsection{Naechste Synthese}

Als naechster Schritt bietet sich eine kurze Ergebnis-Section oder ein
kompaktes Fazit an, das insbesondere folgende Fragen fokussiert:
\begin{itemize}
\item Welche Groesse ist derzeit der beste additive Marker?
\item Welche Groesse ist derzeit der beste multiplikative Marker?
\item Welche Bruecke vermittelt derzeit am besten?
\item Wo sitzen $\gamma$, $-1/12$ und der abc-Anschluss innerhalb des
gemeinsamen Bildes?
\end{itemize}

In dieser Form ist das Material als eigenstaendige, aber gut
kombinierbare LaTeX-Sektion vorbereitet.
 oder % Kombinierbarer Einschub fuer ein groesseres Hauptdokument.
% Vorgesehen fuer % Kombinierbarer Einschub fuer ein groesseres Hauptdokument.
% Vorgesehen fuer % Kombinierbarer Einschub fuer ein groesseres Hauptdokument.
% Vorgesehen fuer \input{Abschluss.tex} oder \include{Abschluss}.

\section{Numerische Unterlegung von Section 7}
\label{sec:section7_numerische_unterlegung}

Die bisherige inhaltliche Struktur von Section~7 laesst sich nun auch
numerisch unterlegen. Ziel ist es, die wesentlichen Marker nicht nur
qualitativ zu benennen, sondern ihre Stabilitaet in zwei verschiedenen
Regimen sowie ihre Kopplung an die Nullstellenseite sichtbar zu machen.

\subsection{Erzeugte Dateien}

Die numerische Auswertung fuehrt auf vier zentrale Ausgabedateien:
\begin{itemize}
\item \texttt{Section 7 -- Bericht},
\item \texttt{Alle Extrapolationswerte},
\item \texttt{Markeruebersicht fuer $s=1$ und $s=-1$},
\item \texttt{Fenstervergleich zur Nullstellenseite}.
\end{itemize}

Diese Dateien bilden zusammen die numerische Basis fuer die folgende
Interpretation.

\subsection{Untersuchte Observablen}

Betrachtet wurden fuenf zentrale Observablen:
\begin{itemize}
\item $K_{\mathrm{add}}$,
\item $G_{P,235}$,
\item $\mathrm{Bridge}_{P,M}$,
\item $\mathrm{Bridge}_{P,W}$,
\item $d_e$ als Abstand zur naechsten E-Primstelle.
\end{itemize}

Fuer jede dieser Groessen wurden zwei komplementaere Regime parallel
untersucht.

\subsection{Das $s=1$-Regime}

Im ersten Regime wurden die Partialsummen
\[
\sum_{p\le X}\frac{F(p)}{p}
\]
gegen
\[
\log\log X+\gamma
\]
gefittet. Die Euler-Mascheroni-Konstante
\[
\gamma \approx 0.5772156649
\]
fungiert dabei nicht nur als kosmetische Verschiebung, sondern als
explizites Zentrum des harmonischen beziehungsweise polnahen
$s=1$-Verhaltens.

\subsection{Das $s=-1$-Regime}

Im zweiten Regime wurden die rohen Summen
\[
\sum_{p\le X} F(p)\,p
\]
mit einem Trendmodell der Form
\[
a\,X\log X+b\,X+c\,\log X+d
\]
approximiert. Anschliessend wurden die Residuen untersucht. Diese
Regularisierung ist zwar noch keine analytische Fortsetzung im strengen
Sinn, sie ist numerisch aber ausreichend stabil, um zwischen ruhigen und
unruhigen Kanaelen zu unterscheiden.

\subsection{Lesart der Markeruebersicht}

Die zentrale Tabelle ist die Markeruebersicht fuer $s=1$ und $s=-1$.
Sie enthaelt fuer jede Observable insbesondere:
\begin{itemize}
\item die Steigung im $\gamma$-zentrierten $s=1$-Fit,
\item den letzten Residualwert,
\item die Gesamtbreite der Residuen,
\item die Koeffizienten des Trendmodells im $s=-1$-Regime.
\end{itemize}

Praktisch ist die Interpretation einfach:
\begin{itemize}
\item Eine kleine \texttt{s1\_residual\_range} bedeutet, dass sich der
Kanal gut an das Bild $\log\log X+\gamma$ anschliesst.
\item Eine kleine \texttt{sminus1\_residual\_range} bedeutet, dass der
Kanal nach Trendabzug im $s=-1$-Regime ruhiger bleibt und damit als
Regularisierungsmarker plausibler wird.
\end{itemize}

\subsection{Vergleich mit der Nullstellenseite}

Dieselben fuenf Observablen wurden zusaetzlich in
$10\,000$er-Fenstern gegen residuale Nullstellenkanaele verglichen:
\begin{itemize}
\item $\Delta_{\mathrm{res}}$,
\item $\Lambda_{\mathrm{res}}$,
\item $\Theta_{Z,\mathrm{res}}$.
\end{itemize}

Die zugehoerige Zusammenfassung liegt im Fenstervergleich zur
Nullstellenseite. Der zentrale Sammelindikator ist
\texttt{score\_abs\_sum}; je groesser dieser Wert ist, desto staerker ist
die kombinierte Kopplung einer Observable an die Nullstellenseite.

Damit liegen erstmals zwei Perspektiven zugleich in derselben
Auswertung vor:
\begin{itemize}
\item Regularisierungsruhe,
\item Nullstellenkopplung.
\end{itemize}

\subsection{Inhaltliches Zwischenurteil}

Der gegenwaertige numerische Stand laesst sich in vier Punkten
zusammenfassen:
\begin{enumerate}
\item Die Euler-Mascheroni-Konstante $\gamma$ gehoert tatsaechlich ins
Zentrum des $s=1$-Regimes. Die Darstellung
\[
\log\log X+\gamma
\]
erweist sich numerisch als sinnvolle Referenzachse.

\item Die additive und die multiplikative Seite trennen sich weiter
klar. $K_{\mathrm{add}}$ bleibt die natuerlichste additive Groesse,
waehrend $G_{P,235}$ weiterhin den staerksten multiplikativen
Haupttraeger bildet.

\item Die Brueckengroessen $\mathrm{Bridge}_{P,M}$ und
$\mathrm{Bridge}_{P,W}$ sind nicht nur dekorative Differenzen, sondern
plausible Kandidaten fuer eine vermittelnde Regularisierung zwischen dem
$\gamma$-zentrierten $s=1$-Bild und dem $s=-1$-regulierten additiven
Bild.

\item Der Abstand $d_e$ zur naechsten E-Primstelle ist inzwischen als
mitlaufende numerische Groesse fest eingebaut und kann als endlicher
Achsenabstieg des e-Anteils gegen den ABC-Kern gelesen werden.
\end{enumerate}

\subsection{Verdichtete Gesamtformulierung}

Die gegenwaertige Gesamtlage laesst sich wie folgt verdichten:

Die Euler-Mascheroni-Konstante $\gamma$ organisiert das harmonische
beziehungsweise polnahe $s=1$-Regime der Vierlingskanaele. Die additive
kleinsche Seite traegt das $s=-1$-Regime, waehrend die glatte
Produktschale die multiplikative Hauptstruktur traegt. Die
Brueckengroessen zwischen glattem und kleinschem Anteil sind die
natuerlichen Kandidaten fuer eine vermittelnde Regularisierung. Der
Abstand zur naechsten E-Primstelle fungiert dabei als endlicher
Achsenabstieg des e-Anteils gegen den ABC-Kern.

\subsection{Naechste Synthese}

Als naechster Schritt bietet sich eine kurze Ergebnis-Section oder ein
kompaktes Fazit an, das insbesondere folgende Fragen fokussiert:
\begin{itemize}
\item Welche Groesse ist derzeit der beste additive Marker?
\item Welche Groesse ist derzeit der beste multiplikative Marker?
\item Welche Bruecke vermittelt derzeit am besten?
\item Wo sitzen $\gamma$, $-1/12$ und der abc-Anschluss innerhalb des
gemeinsamen Bildes?
\end{itemize}

In dieser Form ist das Material als eigenstaendige, aber gut
kombinierbare LaTeX-Sektion vorbereitet.
 oder % Kombinierbarer Einschub fuer ein groesseres Hauptdokument.
% Vorgesehen fuer \input{Abschluss.tex} oder \include{Abschluss}.

\section{Numerische Unterlegung von Section 7}
\label{sec:section7_numerische_unterlegung}

Die bisherige inhaltliche Struktur von Section~7 laesst sich nun auch
numerisch unterlegen. Ziel ist es, die wesentlichen Marker nicht nur
qualitativ zu benennen, sondern ihre Stabilitaet in zwei verschiedenen
Regimen sowie ihre Kopplung an die Nullstellenseite sichtbar zu machen.

\subsection{Erzeugte Dateien}

Die numerische Auswertung fuehrt auf vier zentrale Ausgabedateien:
\begin{itemize}
\item \texttt{Section 7 -- Bericht},
\item \texttt{Alle Extrapolationswerte},
\item \texttt{Markeruebersicht fuer $s=1$ und $s=-1$},
\item \texttt{Fenstervergleich zur Nullstellenseite}.
\end{itemize}

Diese Dateien bilden zusammen die numerische Basis fuer die folgende
Interpretation.

\subsection{Untersuchte Observablen}

Betrachtet wurden fuenf zentrale Observablen:
\begin{itemize}
\item $K_{\mathrm{add}}$,
\item $G_{P,235}$,
\item $\mathrm{Bridge}_{P,M}$,
\item $\mathrm{Bridge}_{P,W}$,
\item $d_e$ als Abstand zur naechsten E-Primstelle.
\end{itemize}

Fuer jede dieser Groessen wurden zwei komplementaere Regime parallel
untersucht.

\subsection{Das $s=1$-Regime}

Im ersten Regime wurden die Partialsummen
\[
\sum_{p\le X}\frac{F(p)}{p}
\]
gegen
\[
\log\log X+\gamma
\]
gefittet. Die Euler-Mascheroni-Konstante
\[
\gamma \approx 0.5772156649
\]
fungiert dabei nicht nur als kosmetische Verschiebung, sondern als
explizites Zentrum des harmonischen beziehungsweise polnahen
$s=1$-Verhaltens.

\subsection{Das $s=-1$-Regime}

Im zweiten Regime wurden die rohen Summen
\[
\sum_{p\le X} F(p)\,p
\]
mit einem Trendmodell der Form
\[
a\,X\log X+b\,X+c\,\log X+d
\]
approximiert. Anschliessend wurden die Residuen untersucht. Diese
Regularisierung ist zwar noch keine analytische Fortsetzung im strengen
Sinn, sie ist numerisch aber ausreichend stabil, um zwischen ruhigen und
unruhigen Kanaelen zu unterscheiden.

\subsection{Lesart der Markeruebersicht}

Die zentrale Tabelle ist die Markeruebersicht fuer $s=1$ und $s=-1$.
Sie enthaelt fuer jede Observable insbesondere:
\begin{itemize}
\item die Steigung im $\gamma$-zentrierten $s=1$-Fit,
\item den letzten Residualwert,
\item die Gesamtbreite der Residuen,
\item die Koeffizienten des Trendmodells im $s=-1$-Regime.
\end{itemize}

Praktisch ist die Interpretation einfach:
\begin{itemize}
\item Eine kleine \texttt{s1\_residual\_range} bedeutet, dass sich der
Kanal gut an das Bild $\log\log X+\gamma$ anschliesst.
\item Eine kleine \texttt{sminus1\_residual\_range} bedeutet, dass der
Kanal nach Trendabzug im $s=-1$-Regime ruhiger bleibt und damit als
Regularisierungsmarker plausibler wird.
\end{itemize}

\subsection{Vergleich mit der Nullstellenseite}

Dieselben fuenf Observablen wurden zusaetzlich in
$10\,000$er-Fenstern gegen residuale Nullstellenkanaele verglichen:
\begin{itemize}
\item $\Delta_{\mathrm{res}}$,
\item $\Lambda_{\mathrm{res}}$,
\item $\Theta_{Z,\mathrm{res}}$.
\end{itemize}

Die zugehoerige Zusammenfassung liegt im Fenstervergleich zur
Nullstellenseite. Der zentrale Sammelindikator ist
\texttt{score\_abs\_sum}; je groesser dieser Wert ist, desto staerker ist
die kombinierte Kopplung einer Observable an die Nullstellenseite.

Damit liegen erstmals zwei Perspektiven zugleich in derselben
Auswertung vor:
\begin{itemize}
\item Regularisierungsruhe,
\item Nullstellenkopplung.
\end{itemize}

\subsection{Inhaltliches Zwischenurteil}

Der gegenwaertige numerische Stand laesst sich in vier Punkten
zusammenfassen:
\begin{enumerate}
\item Die Euler-Mascheroni-Konstante $\gamma$ gehoert tatsaechlich ins
Zentrum des $s=1$-Regimes. Die Darstellung
\[
\log\log X+\gamma
\]
erweist sich numerisch als sinnvolle Referenzachse.

\item Die additive und die multiplikative Seite trennen sich weiter
klar. $K_{\mathrm{add}}$ bleibt die natuerlichste additive Groesse,
waehrend $G_{P,235}$ weiterhin den staerksten multiplikativen
Haupttraeger bildet.

\item Die Brueckengroessen $\mathrm{Bridge}_{P,M}$ und
$\mathrm{Bridge}_{P,W}$ sind nicht nur dekorative Differenzen, sondern
plausible Kandidaten fuer eine vermittelnde Regularisierung zwischen dem
$\gamma$-zentrierten $s=1$-Bild und dem $s=-1$-regulierten additiven
Bild.

\item Der Abstand $d_e$ zur naechsten E-Primstelle ist inzwischen als
mitlaufende numerische Groesse fest eingebaut und kann als endlicher
Achsenabstieg des e-Anteils gegen den ABC-Kern gelesen werden.
\end{enumerate}

\subsection{Verdichtete Gesamtformulierung}

Die gegenwaertige Gesamtlage laesst sich wie folgt verdichten:

Die Euler-Mascheroni-Konstante $\gamma$ organisiert das harmonische
beziehungsweise polnahe $s=1$-Regime der Vierlingskanaele. Die additive
kleinsche Seite traegt das $s=-1$-Regime, waehrend die glatte
Produktschale die multiplikative Hauptstruktur traegt. Die
Brueckengroessen zwischen glattem und kleinschem Anteil sind die
natuerlichen Kandidaten fuer eine vermittelnde Regularisierung. Der
Abstand zur naechsten E-Primstelle fungiert dabei als endlicher
Achsenabstieg des e-Anteils gegen den ABC-Kern.

\subsection{Naechste Synthese}

Als naechster Schritt bietet sich eine kurze Ergebnis-Section oder ein
kompaktes Fazit an, das insbesondere folgende Fragen fokussiert:
\begin{itemize}
\item Welche Groesse ist derzeit der beste additive Marker?
\item Welche Groesse ist derzeit der beste multiplikative Marker?
\item Welche Bruecke vermittelt derzeit am besten?
\item Wo sitzen $\gamma$, $-1/12$ und der abc-Anschluss innerhalb des
gemeinsamen Bildes?
\end{itemize}

In dieser Form ist das Material als eigenstaendige, aber gut
kombinierbare LaTeX-Sektion vorbereitet.
.

\section{Numerische Unterlegung von Section 7}
\label{sec:section7_numerische_unterlegung}

Die bisherige inhaltliche Struktur von Section~7 laesst sich nun auch
numerisch unterlegen. Ziel ist es, die wesentlichen Marker nicht nur
qualitativ zu benennen, sondern ihre Stabilitaet in zwei verschiedenen
Regimen sowie ihre Kopplung an die Nullstellenseite sichtbar zu machen.

\subsection{Erzeugte Dateien}

Die numerische Auswertung fuehrt auf vier zentrale Ausgabedateien:
\begin{itemize}
\item \texttt{Section 7 -- Bericht},
\item \texttt{Alle Extrapolationswerte},
\item \texttt{Markeruebersicht fuer $s=1$ und $s=-1$},
\item \texttt{Fenstervergleich zur Nullstellenseite}.
\end{itemize}

Diese Dateien bilden zusammen die numerische Basis fuer die folgende
Interpretation.

\subsection{Untersuchte Observablen}

Betrachtet wurden fuenf zentrale Observablen:
\begin{itemize}
\item $K_{\mathrm{add}}$,
\item $G_{P,235}$,
\item $\mathrm{Bridge}_{P,M}$,
\item $\mathrm{Bridge}_{P,W}$,
\item $d_e$ als Abstand zur naechsten E-Primstelle.
\end{itemize}

Fuer jede dieser Groessen wurden zwei komplementaere Regime parallel
untersucht.

\subsection{Das $s=1$-Regime}

Im ersten Regime wurden die Partialsummen
\[
\sum_{p\le X}\frac{F(p)}{p}
\]
gegen
\[
\log\log X+\gamma
\]
gefittet. Die Euler-Mascheroni-Konstante
\[
\gamma \approx 0.5772156649
\]
fungiert dabei nicht nur als kosmetische Verschiebung, sondern als
explizites Zentrum des harmonischen beziehungsweise polnahen
$s=1$-Verhaltens.

\subsection{Das $s=-1$-Regime}

Im zweiten Regime wurden die rohen Summen
\[
\sum_{p\le X} F(p)\,p
\]
mit einem Trendmodell der Form
\[
a\,X\log X+b\,X+c\,\log X+d
\]
approximiert. Anschliessend wurden die Residuen untersucht. Diese
Regularisierung ist zwar noch keine analytische Fortsetzung im strengen
Sinn, sie ist numerisch aber ausreichend stabil, um zwischen ruhigen und
unruhigen Kanaelen zu unterscheiden.

\subsection{Lesart der Markeruebersicht}

Die zentrale Tabelle ist die Markeruebersicht fuer $s=1$ und $s=-1$.
Sie enthaelt fuer jede Observable insbesondere:
\begin{itemize}
\item die Steigung im $\gamma$-zentrierten $s=1$-Fit,
\item den letzten Residualwert,
\item die Gesamtbreite der Residuen,
\item die Koeffizienten des Trendmodells im $s=-1$-Regime.
\end{itemize}

Praktisch ist die Interpretation einfach:
\begin{itemize}
\item Eine kleine \texttt{s1\_residual\_range} bedeutet, dass sich der
Kanal gut an das Bild $\log\log X+\gamma$ anschliesst.
\item Eine kleine \texttt{sminus1\_residual\_range} bedeutet, dass der
Kanal nach Trendabzug im $s=-1$-Regime ruhiger bleibt und damit als
Regularisierungsmarker plausibler wird.
\end{itemize}

\subsection{Vergleich mit der Nullstellenseite}

Dieselben fuenf Observablen wurden zusaetzlich in
$10\,000$er-Fenstern gegen residuale Nullstellenkanaele verglichen:
\begin{itemize}
\item $\Delta_{\mathrm{res}}$,
\item $\Lambda_{\mathrm{res}}$,
\item $\Theta_{Z,\mathrm{res}}$.
\end{itemize}

Die zugehoerige Zusammenfassung liegt im Fenstervergleich zur
Nullstellenseite. Der zentrale Sammelindikator ist
\texttt{score\_abs\_sum}; je groesser dieser Wert ist, desto staerker ist
die kombinierte Kopplung einer Observable an die Nullstellenseite.

Damit liegen erstmals zwei Perspektiven zugleich in derselben
Auswertung vor:
\begin{itemize}
\item Regularisierungsruhe,
\item Nullstellenkopplung.
\end{itemize}

\subsection{Inhaltliches Zwischenurteil}

Der gegenwaertige numerische Stand laesst sich in vier Punkten
zusammenfassen:
\begin{enumerate}
\item Die Euler-Mascheroni-Konstante $\gamma$ gehoert tatsaechlich ins
Zentrum des $s=1$-Regimes. Die Darstellung
\[
\log\log X+\gamma
\]
erweist sich numerisch als sinnvolle Referenzachse.

\item Die additive und die multiplikative Seite trennen sich weiter
klar. $K_{\mathrm{add}}$ bleibt die natuerlichste additive Groesse,
waehrend $G_{P,235}$ weiterhin den staerksten multiplikativen
Haupttraeger bildet.

\item Die Brueckengroessen $\mathrm{Bridge}_{P,M}$ und
$\mathrm{Bridge}_{P,W}$ sind nicht nur dekorative Differenzen, sondern
plausible Kandidaten fuer eine vermittelnde Regularisierung zwischen dem
$\gamma$-zentrierten $s=1$-Bild und dem $s=-1$-regulierten additiven
Bild.

\item Der Abstand $d_e$ zur naechsten E-Primstelle ist inzwischen als
mitlaufende numerische Groesse fest eingebaut und kann als endlicher
Achsenabstieg des e-Anteils gegen den ABC-Kern gelesen werden.
\end{enumerate}

\subsection{Verdichtete Gesamtformulierung}

Die gegenwaertige Gesamtlage laesst sich wie folgt verdichten:

Die Euler-Mascheroni-Konstante $\gamma$ organisiert das harmonische
beziehungsweise polnahe $s=1$-Regime der Vierlingskanaele. Die additive
kleinsche Seite traegt das $s=-1$-Regime, waehrend die glatte
Produktschale die multiplikative Hauptstruktur traegt. Die
Brueckengroessen zwischen glattem und kleinschem Anteil sind die
natuerlichen Kandidaten fuer eine vermittelnde Regularisierung. Der
Abstand zur naechsten E-Primstelle fungiert dabei als endlicher
Achsenabstieg des e-Anteils gegen den ABC-Kern.

\subsection{Naechste Synthese}

Als naechster Schritt bietet sich eine kurze Ergebnis-Section oder ein
kompaktes Fazit an, das insbesondere folgende Fragen fokussiert:
\begin{itemize}
\item Welche Groesse ist derzeit der beste additive Marker?
\item Welche Groesse ist derzeit der beste multiplikative Marker?
\item Welche Bruecke vermittelt derzeit am besten?
\item Wo sitzen $\gamma$, $-1/12$ und der abc-Anschluss innerhalb des
gemeinsamen Bildes?
\end{itemize}

In dieser Form ist das Material als eigenstaendige, aber gut
kombinierbare LaTeX-Sektion vorbereitet.
 oder % Kombinierbarer Einschub fuer ein groesseres Hauptdokument.
% Vorgesehen fuer % Kombinierbarer Einschub fuer ein groesseres Hauptdokument.
% Vorgesehen fuer \input{Abschluss.tex} oder \include{Abschluss}.

\section{Numerische Unterlegung von Section 7}
\label{sec:section7_numerische_unterlegung}

Die bisherige inhaltliche Struktur von Section~7 laesst sich nun auch
numerisch unterlegen. Ziel ist es, die wesentlichen Marker nicht nur
qualitativ zu benennen, sondern ihre Stabilitaet in zwei verschiedenen
Regimen sowie ihre Kopplung an die Nullstellenseite sichtbar zu machen.

\subsection{Erzeugte Dateien}

Die numerische Auswertung fuehrt auf vier zentrale Ausgabedateien:
\begin{itemize}
\item \texttt{Section 7 -- Bericht},
\item \texttt{Alle Extrapolationswerte},
\item \texttt{Markeruebersicht fuer $s=1$ und $s=-1$},
\item \texttt{Fenstervergleich zur Nullstellenseite}.
\end{itemize}

Diese Dateien bilden zusammen die numerische Basis fuer die folgende
Interpretation.

\subsection{Untersuchte Observablen}

Betrachtet wurden fuenf zentrale Observablen:
\begin{itemize}
\item $K_{\mathrm{add}}$,
\item $G_{P,235}$,
\item $\mathrm{Bridge}_{P,M}$,
\item $\mathrm{Bridge}_{P,W}$,
\item $d_e$ als Abstand zur naechsten E-Primstelle.
\end{itemize}

Fuer jede dieser Groessen wurden zwei komplementaere Regime parallel
untersucht.

\subsection{Das $s=1$-Regime}

Im ersten Regime wurden die Partialsummen
\[
\sum_{p\le X}\frac{F(p)}{p}
\]
gegen
\[
\log\log X+\gamma
\]
gefittet. Die Euler-Mascheroni-Konstante
\[
\gamma \approx 0.5772156649
\]
fungiert dabei nicht nur als kosmetische Verschiebung, sondern als
explizites Zentrum des harmonischen beziehungsweise polnahen
$s=1$-Verhaltens.

\subsection{Das $s=-1$-Regime}

Im zweiten Regime wurden die rohen Summen
\[
\sum_{p\le X} F(p)\,p
\]
mit einem Trendmodell der Form
\[
a\,X\log X+b\,X+c\,\log X+d
\]
approximiert. Anschliessend wurden die Residuen untersucht. Diese
Regularisierung ist zwar noch keine analytische Fortsetzung im strengen
Sinn, sie ist numerisch aber ausreichend stabil, um zwischen ruhigen und
unruhigen Kanaelen zu unterscheiden.

\subsection{Lesart der Markeruebersicht}

Die zentrale Tabelle ist die Markeruebersicht fuer $s=1$ und $s=-1$.
Sie enthaelt fuer jede Observable insbesondere:
\begin{itemize}
\item die Steigung im $\gamma$-zentrierten $s=1$-Fit,
\item den letzten Residualwert,
\item die Gesamtbreite der Residuen,
\item die Koeffizienten des Trendmodells im $s=-1$-Regime.
\end{itemize}

Praktisch ist die Interpretation einfach:
\begin{itemize}
\item Eine kleine \texttt{s1\_residual\_range} bedeutet, dass sich der
Kanal gut an das Bild $\log\log X+\gamma$ anschliesst.
\item Eine kleine \texttt{sminus1\_residual\_range} bedeutet, dass der
Kanal nach Trendabzug im $s=-1$-Regime ruhiger bleibt und damit als
Regularisierungsmarker plausibler wird.
\end{itemize}

\subsection{Vergleich mit der Nullstellenseite}

Dieselben fuenf Observablen wurden zusaetzlich in
$10\,000$er-Fenstern gegen residuale Nullstellenkanaele verglichen:
\begin{itemize}
\item $\Delta_{\mathrm{res}}$,
\item $\Lambda_{\mathrm{res}}$,
\item $\Theta_{Z,\mathrm{res}}$.
\end{itemize}

Die zugehoerige Zusammenfassung liegt im Fenstervergleich zur
Nullstellenseite. Der zentrale Sammelindikator ist
\texttt{score\_abs\_sum}; je groesser dieser Wert ist, desto staerker ist
die kombinierte Kopplung einer Observable an die Nullstellenseite.

Damit liegen erstmals zwei Perspektiven zugleich in derselben
Auswertung vor:
\begin{itemize}
\item Regularisierungsruhe,
\item Nullstellenkopplung.
\end{itemize}

\subsection{Inhaltliches Zwischenurteil}

Der gegenwaertige numerische Stand laesst sich in vier Punkten
zusammenfassen:
\begin{enumerate}
\item Die Euler-Mascheroni-Konstante $\gamma$ gehoert tatsaechlich ins
Zentrum des $s=1$-Regimes. Die Darstellung
\[
\log\log X+\gamma
\]
erweist sich numerisch als sinnvolle Referenzachse.

\item Die additive und die multiplikative Seite trennen sich weiter
klar. $K_{\mathrm{add}}$ bleibt die natuerlichste additive Groesse,
waehrend $G_{P,235}$ weiterhin den staerksten multiplikativen
Haupttraeger bildet.

\item Die Brueckengroessen $\mathrm{Bridge}_{P,M}$ und
$\mathrm{Bridge}_{P,W}$ sind nicht nur dekorative Differenzen, sondern
plausible Kandidaten fuer eine vermittelnde Regularisierung zwischen dem
$\gamma$-zentrierten $s=1$-Bild und dem $s=-1$-regulierten additiven
Bild.

\item Der Abstand $d_e$ zur naechsten E-Primstelle ist inzwischen als
mitlaufende numerische Groesse fest eingebaut und kann als endlicher
Achsenabstieg des e-Anteils gegen den ABC-Kern gelesen werden.
\end{enumerate}

\subsection{Verdichtete Gesamtformulierung}

Die gegenwaertige Gesamtlage laesst sich wie folgt verdichten:

Die Euler-Mascheroni-Konstante $\gamma$ organisiert das harmonische
beziehungsweise polnahe $s=1$-Regime der Vierlingskanaele. Die additive
kleinsche Seite traegt das $s=-1$-Regime, waehrend die glatte
Produktschale die multiplikative Hauptstruktur traegt. Die
Brueckengroessen zwischen glattem und kleinschem Anteil sind die
natuerlichen Kandidaten fuer eine vermittelnde Regularisierung. Der
Abstand zur naechsten E-Primstelle fungiert dabei als endlicher
Achsenabstieg des e-Anteils gegen den ABC-Kern.

\subsection{Naechste Synthese}

Als naechster Schritt bietet sich eine kurze Ergebnis-Section oder ein
kompaktes Fazit an, das insbesondere folgende Fragen fokussiert:
\begin{itemize}
\item Welche Groesse ist derzeit der beste additive Marker?
\item Welche Groesse ist derzeit der beste multiplikative Marker?
\item Welche Bruecke vermittelt derzeit am besten?
\item Wo sitzen $\gamma$, $-1/12$ und der abc-Anschluss innerhalb des
gemeinsamen Bildes?
\end{itemize}

In dieser Form ist das Material als eigenstaendige, aber gut
kombinierbare LaTeX-Sektion vorbereitet.
 oder % Kombinierbarer Einschub fuer ein groesseres Hauptdokument.
% Vorgesehen fuer \input{Abschluss.tex} oder \include{Abschluss}.

\section{Numerische Unterlegung von Section 7}
\label{sec:section7_numerische_unterlegung}

Die bisherige inhaltliche Struktur von Section~7 laesst sich nun auch
numerisch unterlegen. Ziel ist es, die wesentlichen Marker nicht nur
qualitativ zu benennen, sondern ihre Stabilitaet in zwei verschiedenen
Regimen sowie ihre Kopplung an die Nullstellenseite sichtbar zu machen.

\subsection{Erzeugte Dateien}

Die numerische Auswertung fuehrt auf vier zentrale Ausgabedateien:
\begin{itemize}
\item \texttt{Section 7 -- Bericht},
\item \texttt{Alle Extrapolationswerte},
\item \texttt{Markeruebersicht fuer $s=1$ und $s=-1$},
\item \texttt{Fenstervergleich zur Nullstellenseite}.
\end{itemize}

Diese Dateien bilden zusammen die numerische Basis fuer die folgende
Interpretation.

\subsection{Untersuchte Observablen}

Betrachtet wurden fuenf zentrale Observablen:
\begin{itemize}
\item $K_{\mathrm{add}}$,
\item $G_{P,235}$,
\item $\mathrm{Bridge}_{P,M}$,
\item $\mathrm{Bridge}_{P,W}$,
\item $d_e$ als Abstand zur naechsten E-Primstelle.
\end{itemize}

Fuer jede dieser Groessen wurden zwei komplementaere Regime parallel
untersucht.

\subsection{Das $s=1$-Regime}

Im ersten Regime wurden die Partialsummen
\[
\sum_{p\le X}\frac{F(p)}{p}
\]
gegen
\[
\log\log X+\gamma
\]
gefittet. Die Euler-Mascheroni-Konstante
\[
\gamma \approx 0.5772156649
\]
fungiert dabei nicht nur als kosmetische Verschiebung, sondern als
explizites Zentrum des harmonischen beziehungsweise polnahen
$s=1$-Verhaltens.

\subsection{Das $s=-1$-Regime}

Im zweiten Regime wurden die rohen Summen
\[
\sum_{p\le X} F(p)\,p
\]
mit einem Trendmodell der Form
\[
a\,X\log X+b\,X+c\,\log X+d
\]
approximiert. Anschliessend wurden die Residuen untersucht. Diese
Regularisierung ist zwar noch keine analytische Fortsetzung im strengen
Sinn, sie ist numerisch aber ausreichend stabil, um zwischen ruhigen und
unruhigen Kanaelen zu unterscheiden.

\subsection{Lesart der Markeruebersicht}

Die zentrale Tabelle ist die Markeruebersicht fuer $s=1$ und $s=-1$.
Sie enthaelt fuer jede Observable insbesondere:
\begin{itemize}
\item die Steigung im $\gamma$-zentrierten $s=1$-Fit,
\item den letzten Residualwert,
\item die Gesamtbreite der Residuen,
\item die Koeffizienten des Trendmodells im $s=-1$-Regime.
\end{itemize}

Praktisch ist die Interpretation einfach:
\begin{itemize}
\item Eine kleine \texttt{s1\_residual\_range} bedeutet, dass sich der
Kanal gut an das Bild $\log\log X+\gamma$ anschliesst.
\item Eine kleine \texttt{sminus1\_residual\_range} bedeutet, dass der
Kanal nach Trendabzug im $s=-1$-Regime ruhiger bleibt und damit als
Regularisierungsmarker plausibler wird.
\end{itemize}

\subsection{Vergleich mit der Nullstellenseite}

Dieselben fuenf Observablen wurden zusaetzlich in
$10\,000$er-Fenstern gegen residuale Nullstellenkanaele verglichen:
\begin{itemize}
\item $\Delta_{\mathrm{res}}$,
\item $\Lambda_{\mathrm{res}}$,
\item $\Theta_{Z,\mathrm{res}}$.
\end{itemize}

Die zugehoerige Zusammenfassung liegt im Fenstervergleich zur
Nullstellenseite. Der zentrale Sammelindikator ist
\texttt{score\_abs\_sum}; je groesser dieser Wert ist, desto staerker ist
die kombinierte Kopplung einer Observable an die Nullstellenseite.

Damit liegen erstmals zwei Perspektiven zugleich in derselben
Auswertung vor:
\begin{itemize}
\item Regularisierungsruhe,
\item Nullstellenkopplung.
\end{itemize}

\subsection{Inhaltliches Zwischenurteil}

Der gegenwaertige numerische Stand laesst sich in vier Punkten
zusammenfassen:
\begin{enumerate}
\item Die Euler-Mascheroni-Konstante $\gamma$ gehoert tatsaechlich ins
Zentrum des $s=1$-Regimes. Die Darstellung
\[
\log\log X+\gamma
\]
erweist sich numerisch als sinnvolle Referenzachse.

\item Die additive und die multiplikative Seite trennen sich weiter
klar. $K_{\mathrm{add}}$ bleibt die natuerlichste additive Groesse,
waehrend $G_{P,235}$ weiterhin den staerksten multiplikativen
Haupttraeger bildet.

\item Die Brueckengroessen $\mathrm{Bridge}_{P,M}$ und
$\mathrm{Bridge}_{P,W}$ sind nicht nur dekorative Differenzen, sondern
plausible Kandidaten fuer eine vermittelnde Regularisierung zwischen dem
$\gamma$-zentrierten $s=1$-Bild und dem $s=-1$-regulierten additiven
Bild.

\item Der Abstand $d_e$ zur naechsten E-Primstelle ist inzwischen als
mitlaufende numerische Groesse fest eingebaut und kann als endlicher
Achsenabstieg des e-Anteils gegen den ABC-Kern gelesen werden.
\end{enumerate}

\subsection{Verdichtete Gesamtformulierung}

Die gegenwaertige Gesamtlage laesst sich wie folgt verdichten:

Die Euler-Mascheroni-Konstante $\gamma$ organisiert das harmonische
beziehungsweise polnahe $s=1$-Regime der Vierlingskanaele. Die additive
kleinsche Seite traegt das $s=-1$-Regime, waehrend die glatte
Produktschale die multiplikative Hauptstruktur traegt. Die
Brueckengroessen zwischen glattem und kleinschem Anteil sind die
natuerlichen Kandidaten fuer eine vermittelnde Regularisierung. Der
Abstand zur naechsten E-Primstelle fungiert dabei als endlicher
Achsenabstieg des e-Anteils gegen den ABC-Kern.

\subsection{Naechste Synthese}

Als naechster Schritt bietet sich eine kurze Ergebnis-Section oder ein
kompaktes Fazit an, das insbesondere folgende Fragen fokussiert:
\begin{itemize}
\item Welche Groesse ist derzeit der beste additive Marker?
\item Welche Groesse ist derzeit der beste multiplikative Marker?
\item Welche Bruecke vermittelt derzeit am besten?
\item Wo sitzen $\gamma$, $-1/12$ und der abc-Anschluss innerhalb des
gemeinsamen Bildes?
\end{itemize}

In dieser Form ist das Material als eigenstaendige, aber gut
kombinierbare LaTeX-Sektion vorbereitet.
.

\section{Numerische Unterlegung von Section 7}
\label{sec:section7_numerische_unterlegung}

Die bisherige inhaltliche Struktur von Section~7 laesst sich nun auch
numerisch unterlegen. Ziel ist es, die wesentlichen Marker nicht nur
qualitativ zu benennen, sondern ihre Stabilitaet in zwei verschiedenen
Regimen sowie ihre Kopplung an die Nullstellenseite sichtbar zu machen.

\subsection{Erzeugte Dateien}

Die numerische Auswertung fuehrt auf vier zentrale Ausgabedateien:
\begin{itemize}
\item \texttt{Section 7 -- Bericht},
\item \texttt{Alle Extrapolationswerte},
\item \texttt{Markeruebersicht fuer $s=1$ und $s=-1$},
\item \texttt{Fenstervergleich zur Nullstellenseite}.
\end{itemize}

Diese Dateien bilden zusammen die numerische Basis fuer die folgende
Interpretation.

\subsection{Untersuchte Observablen}

Betrachtet wurden fuenf zentrale Observablen:
\begin{itemize}
\item $K_{\mathrm{add}}$,
\item $G_{P,235}$,
\item $\mathrm{Bridge}_{P,M}$,
\item $\mathrm{Bridge}_{P,W}$,
\item $d_e$ als Abstand zur naechsten E-Primstelle.
\end{itemize}

Fuer jede dieser Groessen wurden zwei komplementaere Regime parallel
untersucht.

\subsection{Das $s=1$-Regime}

Im ersten Regime wurden die Partialsummen
\[
\sum_{p\le X}\frac{F(p)}{p}
\]
gegen
\[
\log\log X+\gamma
\]
gefittet. Die Euler-Mascheroni-Konstante
\[
\gamma \approx 0.5772156649
\]
fungiert dabei nicht nur als kosmetische Verschiebung, sondern als
explizites Zentrum des harmonischen beziehungsweise polnahen
$s=1$-Verhaltens.

\subsection{Das $s=-1$-Regime}

Im zweiten Regime wurden die rohen Summen
\[
\sum_{p\le X} F(p)\,p
\]
mit einem Trendmodell der Form
\[
a\,X\log X+b\,X+c\,\log X+d
\]
approximiert. Anschliessend wurden die Residuen untersucht. Diese
Regularisierung ist zwar noch keine analytische Fortsetzung im strengen
Sinn, sie ist numerisch aber ausreichend stabil, um zwischen ruhigen und
unruhigen Kanaelen zu unterscheiden.

\subsection{Lesart der Markeruebersicht}

Die zentrale Tabelle ist die Markeruebersicht fuer $s=1$ und $s=-1$.
Sie enthaelt fuer jede Observable insbesondere:
\begin{itemize}
\item die Steigung im $\gamma$-zentrierten $s=1$-Fit,
\item den letzten Residualwert,
\item die Gesamtbreite der Residuen,
\item die Koeffizienten des Trendmodells im $s=-1$-Regime.
\end{itemize}

Praktisch ist die Interpretation einfach:
\begin{itemize}
\item Eine kleine \texttt{s1\_residual\_range} bedeutet, dass sich der
Kanal gut an das Bild $\log\log X+\gamma$ anschliesst.
\item Eine kleine \texttt{sminus1\_residual\_range} bedeutet, dass der
Kanal nach Trendabzug im $s=-1$-Regime ruhiger bleibt und damit als
Regularisierungsmarker plausibler wird.
\end{itemize}

\subsection{Vergleich mit der Nullstellenseite}

Dieselben fuenf Observablen wurden zusaetzlich in
$10\,000$er-Fenstern gegen residuale Nullstellenkanaele verglichen:
\begin{itemize}
\item $\Delta_{\mathrm{res}}$,
\item $\Lambda_{\mathrm{res}}$,
\item $\Theta_{Z,\mathrm{res}}$.
\end{itemize}

Die zugehoerige Zusammenfassung liegt im Fenstervergleich zur
Nullstellenseite. Der zentrale Sammelindikator ist
\texttt{score\_abs\_sum}; je groesser dieser Wert ist, desto staerker ist
die kombinierte Kopplung einer Observable an die Nullstellenseite.

Damit liegen erstmals zwei Perspektiven zugleich in derselben
Auswertung vor:
\begin{itemize}
\item Regularisierungsruhe,
\item Nullstellenkopplung.
\end{itemize}

\subsection{Inhaltliches Zwischenurteil}

Der gegenwaertige numerische Stand laesst sich in vier Punkten
zusammenfassen:
\begin{enumerate}
\item Die Euler-Mascheroni-Konstante $\gamma$ gehoert tatsaechlich ins
Zentrum des $s=1$-Regimes. Die Darstellung
\[
\log\log X+\gamma
\]
erweist sich numerisch als sinnvolle Referenzachse.

\item Die additive und die multiplikative Seite trennen sich weiter
klar. $K_{\mathrm{add}}$ bleibt die natuerlichste additive Groesse,
waehrend $G_{P,235}$ weiterhin den staerksten multiplikativen
Haupttraeger bildet.

\item Die Brueckengroessen $\mathrm{Bridge}_{P,M}$ und
$\mathrm{Bridge}_{P,W}$ sind nicht nur dekorative Differenzen, sondern
plausible Kandidaten fuer eine vermittelnde Regularisierung zwischen dem
$\gamma$-zentrierten $s=1$-Bild und dem $s=-1$-regulierten additiven
Bild.

\item Der Abstand $d_e$ zur naechsten E-Primstelle ist inzwischen als
mitlaufende numerische Groesse fest eingebaut und kann als endlicher
Achsenabstieg des e-Anteils gegen den ABC-Kern gelesen werden.
\end{enumerate}

\subsection{Verdichtete Gesamtformulierung}

Die gegenwaertige Gesamtlage laesst sich wie folgt verdichten:

Die Euler-Mascheroni-Konstante $\gamma$ organisiert das harmonische
beziehungsweise polnahe $s=1$-Regime der Vierlingskanaele. Die additive
kleinsche Seite traegt das $s=-1$-Regime, waehrend die glatte
Produktschale die multiplikative Hauptstruktur traegt. Die
Brueckengroessen zwischen glattem und kleinschem Anteil sind die
natuerlichen Kandidaten fuer eine vermittelnde Regularisierung. Der
Abstand zur naechsten E-Primstelle fungiert dabei als endlicher
Achsenabstieg des e-Anteils gegen den ABC-Kern.

\subsection{Naechste Synthese}

Als naechster Schritt bietet sich eine kurze Ergebnis-Section oder ein
kompaktes Fazit an, das insbesondere folgende Fragen fokussiert:
\begin{itemize}
\item Welche Groesse ist derzeit der beste additive Marker?
\item Welche Groesse ist derzeit der beste multiplikative Marker?
\item Welche Bruecke vermittelt derzeit am besten?
\item Wo sitzen $\gamma$, $-1/12$ und der abc-Anschluss innerhalb des
gemeinsamen Bildes?
\end{itemize}

In dieser Form ist das Material als eigenstaendige, aber gut
kombinierbare LaTeX-Sektion vorbereitet.
.

\section{Numerische Unterlegung von Section 7}
\label{sec:section7_numerische_unterlegung}

Die bisherige inhaltliche Struktur von Section~7 laesst sich nun auch
numerisch unterlegen. Ziel ist es, die wesentlichen Marker nicht nur
qualitativ zu benennen, sondern ihre Stabilitaet in zwei verschiedenen
Regimen sowie ihre Kopplung an die Nullstellenseite sichtbar zu machen.

\subsection{Erzeugte Dateien}

Die numerische Auswertung fuehrt auf vier zentrale Ausgabedateien:
\begin{itemize}
\item \texttt{Section 7 -- Bericht},
\item \texttt{Alle Extrapolationswerte},
\item \texttt{Markeruebersicht fuer $s=1$ und $s=-1$},
\item \texttt{Fenstervergleich zur Nullstellenseite}.
\end{itemize}

Diese Dateien bilden zusammen die numerische Basis fuer die folgende
Interpretation.

\subsection{Untersuchte Observablen}

Betrachtet wurden fuenf zentrale Observablen:
\begin{itemize}
\item $K_{\mathrm{add}}$,
\item $G_{P,235}$,
\item $\mathrm{Bridge}_{P,M}$,
\item $\mathrm{Bridge}_{P,W}$,
\item $d_e$ als Abstand zur naechsten E-Primstelle.
\end{itemize}

Fuer jede dieser Groessen wurden zwei komplementaere Regime parallel
untersucht.

\subsection{Das $s=1$-Regime}

Im ersten Regime wurden die Partialsummen
\[
\sum_{p\le X}\frac{F(p)}{p}
\]
gegen
\[
\log\log X+\gamma
\]
gefittet. Die Euler-Mascheroni-Konstante
\[
\gamma \approx 0.5772156649
\]
fungiert dabei nicht nur als kosmetische Verschiebung, sondern als
explizites Zentrum des harmonischen beziehungsweise polnahen
$s=1$-Verhaltens.

\subsection{Das $s=-1$-Regime}

Im zweiten Regime wurden die rohen Summen
\[
\sum_{p\le X} F(p)\,p
\]
mit einem Trendmodell der Form
\[
a\,X\log X+b\,X+c\,\log X+d
\]
approximiert. Anschliessend wurden die Residuen untersucht. Diese
Regularisierung ist zwar noch keine analytische Fortsetzung im strengen
Sinn, sie ist numerisch aber ausreichend stabil, um zwischen ruhigen und
unruhigen Kanaelen zu unterscheiden.

\subsection{Lesart der Markeruebersicht}

Die zentrale Tabelle ist die Markeruebersicht fuer $s=1$ und $s=-1$.
Sie enthaelt fuer jede Observable insbesondere:
\begin{itemize}
\item die Steigung im $\gamma$-zentrierten $s=1$-Fit,
\item den letzten Residualwert,
\item die Gesamtbreite der Residuen,
\item die Koeffizienten des Trendmodells im $s=-1$-Regime.
\end{itemize}

Praktisch ist die Interpretation einfach:
\begin{itemize}
\item Eine kleine \texttt{s1\_residual\_range} bedeutet, dass sich der
Kanal gut an das Bild $\log\log X+\gamma$ anschliesst.
\item Eine kleine \texttt{sminus1\_residual\_range} bedeutet, dass der
Kanal nach Trendabzug im $s=-1$-Regime ruhiger bleibt und damit als
Regularisierungsmarker plausibler wird.
\end{itemize}

\subsection{Vergleich mit der Nullstellenseite}

Dieselben fuenf Observablen wurden zusaetzlich in
$10\,000$er-Fenstern gegen residuale Nullstellenkanaele verglichen:
\begin{itemize}
\item $\Delta_{\mathrm{res}}$,
\item $\Lambda_{\mathrm{res}}$,
\item $\Theta_{Z,\mathrm{res}}$.
\end{itemize}

Die zugehoerige Zusammenfassung liegt im Fenstervergleich zur
Nullstellenseite. Der zentrale Sammelindikator ist
\texttt{score\_abs\_sum}; je groesser dieser Wert ist, desto staerker ist
die kombinierte Kopplung einer Observable an die Nullstellenseite.

Damit liegen erstmals zwei Perspektiven zugleich in derselben
Auswertung vor:
\begin{itemize}
\item Regularisierungsruhe,
\item Nullstellenkopplung.
\end{itemize}

\subsection{Inhaltliches Zwischenurteil}

Der gegenwaertige numerische Stand laesst sich in vier Punkten
zusammenfassen:
\begin{enumerate}
\item Die Euler-Mascheroni-Konstante $\gamma$ gehoert tatsaechlich ins
Zentrum des $s=1$-Regimes. Die Darstellung
\[
\log\log X+\gamma
\]
erweist sich numerisch als sinnvolle Referenzachse.

\item Die additive und die multiplikative Seite trennen sich weiter
klar. $K_{\mathrm{add}}$ bleibt die natuerlichste additive Groesse,
waehrend $G_{P,235}$ weiterhin den staerksten multiplikativen
Haupttraeger bildet.

\item Die Brueckengroessen $\mathrm{Bridge}_{P,M}$ und
$\mathrm{Bridge}_{P,W}$ sind nicht nur dekorative Differenzen, sondern
plausible Kandidaten fuer eine vermittelnde Regularisierung zwischen dem
$\gamma$-zentrierten $s=1$-Bild und dem $s=-1$-regulierten additiven
Bild.

\item Der Abstand $d_e$ zur naechsten E-Primstelle ist inzwischen als
mitlaufende numerische Groesse fest eingebaut und kann als endlicher
Achsenabstieg des e-Anteils gegen den ABC-Kern gelesen werden.
\end{enumerate}

\subsection{Verdichtete Gesamtformulierung}

Die gegenwaertige Gesamtlage laesst sich wie folgt verdichten:

Die Euler-Mascheroni-Konstante $\gamma$ organisiert das harmonische
beziehungsweise polnahe $s=1$-Regime der Vierlingskanaele. Die additive
kleinsche Seite traegt das $s=-1$-Regime, waehrend die glatte
Produktschale die multiplikative Hauptstruktur traegt. Die
Brueckengroessen zwischen glattem und kleinschem Anteil sind die
natuerlichen Kandidaten fuer eine vermittelnde Regularisierung. Der
Abstand zur naechsten E-Primstelle fungiert dabei als endlicher
Achsenabstieg des e-Anteils gegen den ABC-Kern.

\subsection{Naechste Synthese}

Als naechster Schritt bietet sich eine kurze Ergebnis-Section oder ein
kompaktes Fazit an, das insbesondere folgende Fragen fokussiert:
\begin{itemize}
\item Welche Groesse ist derzeit der beste additive Marker?
\item Welche Groesse ist derzeit der beste multiplikative Marker?
\item Welche Bruecke vermittelt derzeit am besten?
\item Wo sitzen $\gamma$, $-1/12$ und der abc-Anschluss innerhalb des
gemeinsamen Bildes?
\end{itemize}

In dieser Form ist das Material als eigenstaendige, aber gut
kombinierbare LaTeX-Sektion vorbereitet.
.

\section{Numerische Unterlegung von Section 7}
\label{sec:section7_numerische_unterlegung}

Die bisherige inhaltliche Struktur von Section~7 laesst sich nun auch
numerisch unterlegen. Ziel ist es, die wesentlichen Marker nicht nur
qualitativ zu benennen, sondern ihre Stabilitaet in zwei verschiedenen
Regimen sowie ihre Kopplung an die Nullstellenseite sichtbar zu machen.

\subsection{Erzeugte Dateien}

Die numerische Auswertung fuehrt auf vier zentrale Ausgabedateien:
\begin{itemize}
\item \texttt{Section 7 -- Bericht},
\item \texttt{Alle Extrapolationswerte},
\item \texttt{Markeruebersicht fuer $s=1$ und $s=-1$},
\item \texttt{Fenstervergleich zur Nullstellenseite}.
\end{itemize}

Diese Dateien bilden zusammen die numerische Basis fuer die folgende
Interpretation.

\subsection{Untersuchte Observablen}

Betrachtet wurden fuenf zentrale Observablen:
\begin{itemize}
\item $K_{\mathrm{add}}$,
\item $G_{P,235}$,
\item $\mathrm{Bridge}_{P,M}$,
\item $\mathrm{Bridge}_{P,W}$,
\item $d_e$ als Abstand zur naechsten E-Primstelle.
\end{itemize}

Fuer jede dieser Groessen wurden zwei komplementaere Regime parallel
untersucht.

\subsection{Das $s=1$-Regime}

Im ersten Regime wurden die Partialsummen
\[
\sum_{p\le X}\frac{F(p)}{p}
\]
gegen
\[
\log\log X+\gamma
\]
gefittet. Die Euler-Mascheroni-Konstante
\[
\gamma \approx 0.5772156649
\]
fungiert dabei nicht nur als kosmetische Verschiebung, sondern als
explizites Zentrum des harmonischen beziehungsweise polnahen
$s=1$-Verhaltens.

\subsection{Das $s=-1$-Regime}

Im zweiten Regime wurden die rohen Summen
\[
\sum_{p\le X} F(p)\,p
\]
mit einem Trendmodell der Form
\[
a\,X\log X+b\,X+c\,\log X+d
\]
approximiert. Anschliessend wurden die Residuen untersucht. Diese
Regularisierung ist zwar noch keine analytische Fortsetzung im strengen
Sinn, sie ist numerisch aber ausreichend stabil, um zwischen ruhigen und
unruhigen Kanaelen zu unterscheiden.

\subsection{Lesart der Markeruebersicht}

Die zentrale Tabelle ist die Markeruebersicht fuer $s=1$ und $s=-1$.
Sie enthaelt fuer jede Observable insbesondere:
\begin{itemize}
\item die Steigung im $\gamma$-zentrierten $s=1$-Fit,
\item den letzten Residualwert,
\item die Gesamtbreite der Residuen,
\item die Koeffizienten des Trendmodells im $s=-1$-Regime.
\end{itemize}

Praktisch ist die Interpretation einfach:
\begin{itemize}
\item Eine kleine \texttt{s1\_residual\_range} bedeutet, dass sich der
Kanal gut an das Bild $\log\log X+\gamma$ anschliesst.
\item Eine kleine \texttt{sminus1\_residual\_range} bedeutet, dass der
Kanal nach Trendabzug im $s=-1$-Regime ruhiger bleibt und damit als
Regularisierungsmarker plausibler wird.
\end{itemize}

\subsection{Vergleich mit der Nullstellenseite}

Dieselben fuenf Observablen wurden zusaetzlich in
$10\,000$er-Fenstern gegen residuale Nullstellenkanaele verglichen:
\begin{itemize}
\item $\Delta_{\mathrm{res}}$,
\item $\Lambda_{\mathrm{res}}$,
\item $\Theta_{Z,\mathrm{res}}$.
\end{itemize}

Die zugehoerige Zusammenfassung liegt im Fenstervergleich zur
Nullstellenseite. Der zentrale Sammelindikator ist
\texttt{score\_abs\_sum}; je groesser dieser Wert ist, desto staerker ist
die kombinierte Kopplung einer Observable an die Nullstellenseite.

Damit liegen erstmals zwei Perspektiven zugleich in derselben
Auswertung vor:
\begin{itemize}
\item Regularisierungsruhe,
\item Nullstellenkopplung.
\end{itemize}

\subsection{Inhaltliches Zwischenurteil}

Der gegenwaertige numerische Stand laesst sich in vier Punkten
zusammenfassen:
\begin{enumerate}
\item Die Euler-Mascheroni-Konstante $\gamma$ gehoert tatsaechlich ins
Zentrum des $s=1$-Regimes. Die Darstellung
\[
\log\log X+\gamma
\]
erweist sich numerisch als sinnvolle Referenzachse.

\item Die additive und die multiplikative Seite trennen sich weiter
klar. $K_{\mathrm{add}}$ bleibt die natuerlichste additive Groesse,
waehrend $G_{P,235}$ weiterhin den staerksten multiplikativen
Haupttraeger bildet.

\item Die Brueckengroessen $\mathrm{Bridge}_{P,M}$ und
$\mathrm{Bridge}_{P,W}$ sind nicht nur dekorative Differenzen, sondern
plausible Kandidaten fuer eine vermittelnde Regularisierung zwischen dem
$\gamma$-zentrierten $s=1$-Bild und dem $s=-1$-regulierten additiven
Bild.

\item Der Abstand $d_e$ zur naechsten E-Primstelle ist inzwischen als
mitlaufende numerische Groesse fest eingebaut und kann als endlicher
Achsenabstieg des e-Anteils gegen den ABC-Kern gelesen werden.
\end{enumerate}

\subsection{Verdichtete Gesamtformulierung}

Die gegenwaertige Gesamtlage laesst sich wie folgt verdichten:

Die Euler-Mascheroni-Konstante $\gamma$ organisiert das harmonische
beziehungsweise polnahe $s=1$-Regime der Vierlingskanaele. Die additive
kleinsche Seite traegt das $s=-1$-Regime, waehrend die glatte
Produktschale die multiplikative Hauptstruktur traegt. Die
Brueckengroessen zwischen glattem und kleinschem Anteil sind die
natuerlichen Kandidaten fuer eine vermittelnde Regularisierung. Der
Abstand zur naechsten E-Primstelle fungiert dabei als endlicher
Achsenabstieg des e-Anteils gegen den ABC-Kern.

\subsection{Naechste Synthese}

Als naechster Schritt bietet sich eine kurze Ergebnis-Section oder ein
kompaktes Fazit an, das insbesondere folgende Fragen fokussiert:
\begin{itemize}
\item Welche Groesse ist derzeit der beste additive Marker?
\item Welche Groesse ist derzeit der beste multiplikative Marker?
\item Welche Bruecke vermittelt derzeit am besten?
\item Wo sitzen $\gamma$, $-1/12$ und der abc-Anschluss innerhalb des
gemeinsamen Bildes?
\end{itemize}

In dieser Form ist das Material als eigenstaendige, aber gut
kombinierbare LaTeX-Sektion vorbereitet.
